\section{Shared results V: divisors, Iwasawa decompositions and interpolation}\label{sec:pf:sh-div}

\providecommand{\shdivGBig}{\sfname{GBig}}
\providecommand{\shdivGC}{\Gamma_{\C}}
\providecommand{\shdivEm}{\mathrm{E}_{\mathrm{main}}}
\providecommand{\shdivLi}{\operatorname{Li}}
\providecommand{\shdivPic}{\operatorname{Pic}}
\providecommand{\shdivex}{\operatorname{ex}}
\providecommand{\shdivEv}{\sfname{Ev}}

This section proves nineteen results of the formal development \cite{LeanSource2026} which are used by several of the
preceding sections of Part~II. They are of very different nature and size. Some are classical theorems (or formalized
forms of them): the boundedness of completed Hecke $L$-functions in vertical strips, the prime number theorem for arithmetic
progressions with an exceptional zero, the finiteness of the volume of an adelic unitary quotient. For these we state
precisely what the formal statement asserts, compare it with the published statement and describe the formal proof. Others
are re-packagings or specialisations which get a short proof. The remaining ones (rationality and integrality of theta series,
Galois conjugates of rational theta lifts, the volume ratio identity, the cut inequality for Weyl modules, the class sum span)
rest on long formal proofs; for them we give the statement, the complete deduction from the main ingredients, and the
statements of the ingredients, the proofs of the deepest of which are not reproduced here (they are listed with their
content where they are used).

The section is organised by topic.
\begin{itemize}
\item \Cref{sec:pf:sh-div:hecke}: Hecke $L$-functions of CM fields (used in \cref{sec:pf:fit,sec:pf:rowprops,sec:pf:liftinj}):
boundedness in strips (\cref{pf:sh-div:strips}), the absolute value form of the functional equation
(\cref{pf:sh-div:feabs}), non-vanishing at integers (\cref{pf:sh-div:lne}), the field data of a host
(\cref{pf:sh-div:fielddat}) and the Galois conjugates of toric ratios (\cref{pf:sh-div:toric}).
\item \Cref{sec:pf:sh-div:primes}: primes in arithmetic progressions and split primes (used in
\cref{sec:pf:stageU,sec:pf:top-supply}): \cref{pf:sh-div:primesrel,pf:sh-div:apinterval,pf:sh-div:split}.
\item \Cref{sec:pf:sh-div:vol}: volumes (used in \cref{sec:pf:fs,sec:pf:stageU}): finiteness of the volume of the
adelic unitary quotient (\cref{pf:sh-div:finvol}) and the volume ratio identity from the one-vector Siegel formula
(\cref{pf:sh-div:volratio}).
\item \Cref{sec:pf:sh-div:theta}: theta series (used in \cref{sec:pf:liftinj,sec:pf:onto3,sec:pf:rowprops,sec:pf:stageU,%
sec:pf:open0,sec:pf:rows}): rationality and integrality (\cref{pf:sh-div:series}), conjugates of rational theta lifts
(\cref{pf:sh-div:conj}), the class sum span (\cref{pf:sh-div:classsum}), the rational row test
(\cref{pf:sh-div:rowrat}) and symmetric column forms (\cref{pf:sh-div:symcol}).
\item \Cref{sec:pf:sh-div:frames}: frames, twisted places, numerical rows and the cut inequality (used in
\cref{sec:pf:baselat1,sec:pf:baselat2,sec:pf:bundles,sec:pf:open0,sec:pf:stageU,sec:pf:onto2}):
\cref{pf:sh-div:hostframes,pf:sh-div:twisted,pf:sh-div:szrow,pf:sh-div:excut}.
\end{itemize}
Throughout, the objects of the development (global data $D$, hosts, real data, faces, grids, stages, contexts) are those of
\cref{sec:hosts,sec:families} and of the sections of Part~II quoted at each place; we recall what is needed.

%==========================================================================================================================
\subsection{Hecke \texorpdfstring{$L$}{L}-functions of CM fields}\label{sec:pf:sh-div:hecke}

\paragraph{Conventions.} Let $E$ be a number field with idele group $\A_E^\times$. A \emph{Hecke character} of $E$ is a
continuous homomorphism $\chi:\A_E^\times\to\C^\times$ trivial on the diagonal image of $E^\times$; it is \emph{unitary} if
$|\chi|=1$. For a finite place $v$ let $\chi_v$ be the restriction of $\chi$ to $E_v^\times$, fix $\varpi_v\in E^\times$ which is a
uniformiser at $v$, and put $\chi(v)=\chi_v(\varpi_v)$. For $k\in E^\times$ let $\chi_\infty(k)$ be the value of $\chi$ at the idele
which is $k$ at the archimedean places and $1$ elsewhere. The character $\chi$ is \emph{unramified} at $v$ if $\chi_v$ is trivial
on $\OO_v^\times$, and $a\in\N$ is the \emph{conductor exponent} of $\chi$ at $v$ if for every $m\in\N$, $\chi_v$ is trivial on
the units $x\equiv1\bmod\mathfrak p_v^m$ (all units if $m=0$) exactly when $a\le m$; we then write $a=a_v(\chi)$.

Let $E$ be totally complex and let $\varphi=(\varphi_i)_{i\in I}$ be a finite family of embeddings $E\to\C$. It is a
\emph{CM type} if every embedding $\sigma:E\to\C$ is equal to $\varphi_i$ or to $\overline{\varphi_i}$ for exactly one $i$. A Hecke
character $\chi$ has \emph{infinity type} $n\in\Z^I$ (with respect to $\varphi$) if
\begin{equation}\label{eq:pf:sh-div:inftype}
\chi_\infty(k)=\prod_{i\in I}\Bigl(\frac{\varphi_i(k)}{|\varphi_i(k)|}\Bigr)^{n_i}\qquad(k\in E^\times),
\end{equation}
and \emph{weight} $m\in\Z^I$ if it has infinity type $2m$, i.e.\ $\chi_\infty(k)=\prod_i(\varphi_i(k)/\overline{\varphi_i(k)})^{m_i}$.
The Euler factor at a finite place $v$ is $L_v(s,\chi)=(1-\chi(v)\Nm v^{-s})^{-1}$ if $\chi$ is unramified at $v$ and $1$ otherwise,
and for $n\in\Z^I$ we put
\[
\Gamma_\infty(n,s)=\prod_{i\in I}\shdivGC\bigl(s+\tfrac{|n_i|}2\bigr),\qquad \shdivGC(s)=2(2\pi)^{-s}\Gamma(s).
\]
A \emph{completed $L$-function} of $(\chi,n)$ is an entire function $\Lambda$ such that for $\Re s>1$ the Euler product
$\prod_vL_v(s,\chi)$ converges unconditionally to a number $L$ with $\Lambda(s)=\Gamma_\infty(n,s)L$. A \emph{Hecke $L$-function} of $\chi$ is
an entire function $\Lambda$ equal to the Euler product $\prod_vL_v(s,\chi)$ for $\Re s>1$ (no archimedean factor). Two completed
$L$-functions (or two Hecke $L$-functions) of the same data coincide, since they agree on $\Re s>1$. A function $\Lambda$ is
\emph{bounded in vertical strips} if for all real $a<b$ it is bounded on $\{a\le\Re s\le b\}$.

For a unitary $\chi$ and $\Re s>1$ the Euler product converges absolutely and is the exponential of
$\sum_v-\log(1-c_v\Nm v^{-s})$, where $c_v=\chi(v)$ if $\chi$ is unramified at $v$ and $c_v=0$ otherwise
(\lref{GlobalTheory.HeckeOne.hasProd_exp}); in particular it does not vanish there.

\begin{theorem}[Completed Hecke $L$-functions are bounded in vertical strips]\label{pf:sh-div:strips}
Let $E$ be a totally complex number field, $\varphi$ a CM type of $E$ indexed by $I$, and $\chi$ a unitary Hecke character of $E$
of infinity type $n\in\Z^I$ with $n\ne0$. Then $(\chi,n)$ has a completed $L$-function $\Lambda$ which is bounded in vertical
strips.
\formallabel{GlobalTheory.HeckeStrip.completedLBoundedInStrips}
\end{theorem}

This is Hecke's theorem on the analytic continuation of the $L$-series of Gr\"o\ss{}encharaktere
\cite{Hecke1920}, in Tate's adelic form \cite{Tate1967thesis}; see \cite[Ch.~VII, \S8]{Neukirch1999}. Since $n\ne0$, $\chi$ is
not of the form $|\cdot|^{it}$, so the completed $L$-function has no poles; boundedness in vertical strips follows from the
representation of $\Lambda$ as a Mellin transform of a theta series, split at $t=1$, which is uniformly convergent on vertical
strips. The formal statement differs from the published one in two harmless ways: the factor $(|d_E|\Nm\mathfrak f(\chi))^{s/2}$ is
omitted from $\Lambda$ (it is entire, non-vanishing and bounded on vertical strips), and only the existence of an entire
function with the stated Euler product and the bound on strips is asserted. The formal proof follows Hecke's argument.

\begin{proof}
\emph{Step 1: the ideal character.} Since $\chi$ is unitary, there are a non-zero ideal $\mathfrak m$ of $\OO_E$ and a
multiplicative function $X$ on the ideals of $\OO_E$ with $X(0)=0$, $X(\OO_E)=1$, $|X(\mathfrak a)|=1$ for non-zero $\mathfrak a$
prime to $\mathfrak m$, $X(\mathfrak a)=0$ if $\mathfrak a$ is not prime to $\mathfrak m$, $X(\mathfrak p_v)=\chi(v)$ if $\chi$ is
unramified at $v$ and $0$ otherwise, and $X(\beta\OO_E)=\chi_\infty(\beta)^{-1}$ for $\beta\in\OO_E$ with $\beta\equiv1\bmod\mathfrak m$
(\lref{GlobalTheory.T.exists_idealCharacter_of_unitary_heckeChar}). By \eqref{eq:pf:sh-div:inftype},
$X(\beta\OO_E)=w(\beta)$ for such $\beta\ne0$, where $w(\alpha)=\prod_i(\varphi_i(\alpha)/|\varphi_i(\alpha)|)^{-n_i}$.

\emph{Step 2: decomposition into partial series.} For a function $F$ on $\OO_E$ put
$P_F(s)=\sum_{(\alpha)}F(\alpha)w(\alpha)\Nm(\alpha)^{-s}$, the sum over the non-zero principal ideals $(\alpha)$; it is well defined
when $F(u\alpha)w(u\alpha)=F(\alpha)w(\alpha)$ for all units $u$. Grouping the ideals prime to $\mathfrak m$ by their ray classes and
multiplying each class by a fixed integral ideal of the inverse class, one obtains finitely many constants $c_k\in\C$, integers
$B_k\ge1$, non-zero ideals $M_k$ and functions $F_k$ on $\OO_E$, periodic modulo $M_k$ and satisfying the unit condition, with
\[
\sum_{\mathfrak a}X(\mathfrak a)\Nm\mathfrak a^{-s}=\sum_kc_kB_k^{\,s}P_{F_k}(s)\qquad(\Re s>1)
\]
(\lref{GlobalTheory.T.heckeLSeries_eq_sum_principalPartialSeries}).

\emph{Step 3: Hecke's theta argument.} For every such $F$ (periodic modulo a non-zero ideal, with the unit condition) the
function $\Gamma_\infty(n,s)P_F(s)$, $\Re s>1$, extends to an entire function bounded in vertical strips
(\lref{GlobalTheory.T.completed_heckePartialLSeries_entire_bounded_strips}). The formal proof is Hecke's: for the theta
function $\theta_F(y)$, $y\in\R^I$, attached to $F$ and $n$ (a sum over $\alpha\in\OO_E$ weighted by $F(\alpha)$, the coordinates
$\varphi_i(\alpha)$ and $n$), Poisson summation gives $\theta_F(y)=\kappa e^{-\sum_iy_i}\theta_a(-y)$ with a dual coefficient function $a$
whose non-zero arguments stay away from $0$ (\lref{GlobalTheory.T.heckeThetaKernel_poisson_summation}); $\Gamma_\infty(n,s)P_F(s)$ is
written as an integral of $\theta_F$ over a fundamental domain for the action of the units (obtained from a measurable partition of
unity by unit-logarithm translates); the integral is split by the functional equation of the Mellin transform into two integrals
of Laplace type against rapidly decreasing functions, which are entire and bounded in vertical strips
(\lref{GlobalTheory.T.mellin_integral_split_functional_equation} and \lref{GlobalTheory.T.laplace_integral_entire_boundedInVerticalStrips}).

\emph{Step 4: conclusion.} Let $\Lambda_k$ be the entire functions of Step~3 for $F_k$ and $\Lambda(s)=\sum_kc_kB_k^{\,s}\Lambda_k(s)$.
It is entire. For $\Re s>1$, the Dirichlet series $\sum X(\mathfrak a)\Nm\mathfrak a^{-s}$ equals the exponential of
$\sum_v-\log(1-X(\mathfrak p_v)\Nm v^{-s})$ (\lref{GlobalTheory.tsum_ideal_mul_absNorm_cpow_eq_exp_tsum}: unique factorisation and
absolute convergence for a completely multiplicative $X$ with $|X|\le1$), which by Step~1 and
\lref{GlobalTheory.HeckeOne.hasProd_exp} is the limit of the Euler product of $\chi$; by Step~2,
$\Lambda(s)=\Gamma_\infty(n,s)\prod_vL_v(s,\chi)$. Finally on $a\le\Re s\le b$ we have $|B_k^{\,s}|\le B_k^{\,b}$ and $|\Lambda_k|\le C_k$, so
$|\Lambda|\le\sum_k|c_k|B_k^{\,b}C_k$.
\end{proof}

\begin{definition}[Root numbers]\label{pf:sh-div:rootnumbers}
Let $E$ be a number field, $\chi$ a Hecke character of $E$ and $a=(a_v)$ a finitely supported family of natural numbers indexed by
the finite places. Let $d_v$ be the exponent of $\mathfrak p_v$ in the different of $E/\Q$. For $c\in E^\times$ let $G_v(\chi,a_v,c)$
be the local Gauss sum of $\chi$ at $v$ modulo $\mathfrak p_v^{a_v}$ attached to $c$ \cite{LeanSource2026}. Put
\[
\varepsilon_v(\chi,a_v,s)=\chi_v(\varpi_v^{a_v+d_v})\Nm v^{-s(a_v+d_v)}\Nm v^{d_v/2}\cdot
\begin{cases}1,&a_v=0,\\ G_v(\chi,a_v,\varpi_v^{a_v+d_v}),&a_v\ge1,\end{cases}
\]
$\varepsilon_f(\chi,a,s)=\prod_v\varepsilon_v(\chi,a_v,s)$ (a finite product: $\varepsilon_v=1$ if $a_v=d_v=0$) and, for $n\in\Z^I$,
$\varepsilon_\infty(n)=(-i)^{\sum_i|n_i|}$.
\end{definition}

The \emph{functional equation of Hecke $L$-functions} $\sfname{HeckeTateFunctionalEquation}$ is the following statement (a
hypothesis of the development, labelled as a published result, \cite{Hecke1920,Tate1967thesis}): for every totally complex $E$ with
a CM type $\varphi$, every unitary Hecke character $\chi$ of infinity type $n\ne0$ and conductor exponents $a=(a_v(\chi))_v$, there
are completed $L$-functions $\Lambda$ of $(\chi,n)$ and $\Lambda'$ of $(\chi^{-1},-n)$ with
\begin{equation}\label{eq:pf:sh-div:fe}
\Lambda(s)=\varepsilon_\infty(n)\,\varepsilon_f(\chi,a,s)\,\Lambda'(1-s)\qquad(s\in\C).
\end{equation}
Its \emph{absolute value form} $\sfname{HeckeFunctionalEquationAbs}$ is the statement: for $E,\varphi,\chi,n,a$ as above and
\emph{all} completed $L$-functions $\Lambda$ of $(\chi,n)$ and $\Lambda'$ of $(\chi^{-1},-n)$ there is $\varepsilon\in\C$ with
$|\varepsilon|=1$ and
\begin{equation}\label{eq:pf:sh-div:feabs}
\Lambda(s)=\varepsilon\,\Bigl(|d_E|\prod_v\Nm v^{a_v}\Bigr)^{1/2-s}\Lambda'(1-s)\qquad(s\in\C).
\end{equation}

\begin{proposition}[The functional equation in absolute value form]\label{pf:sh-div:feabs}
The functional equation of Hecke $L$-functions, $\sfname{HeckeTateFunctionalEquation}$, implies its absolute value form,
$\sfname{HeckeFunctionalEquationAbs}$.
\formallabel{GlobalTheory.Published.HeckeTateFunctionalEquation.abs}
\formallabel{GlobalTheory.HeckeEpsAbs.locEps_eq_mul}\formallabel{GlobalTheory.HeckeEpsAbs.finEps_eq}
\formallabel{GlobalTheory.HeckeEpsAbs.norm_locEps_half}\formallabel{GlobalTheory.HeckeEpsAbs.norm_finEps_half}
\formallabel{GlobalTheory.HeckeEpsAbs.norm_archEps}
\end{proposition}

\begin{proof}
Let $\Lambda_0,\Lambda'_0$ be the functions given by \eqref{eq:pf:sh-div:fe}. By uniqueness of completed $L$-functions,
$\Lambda=\Lambda_0$ and $\Lambda'=\Lambda'_0$. Put $\varepsilon=\varepsilon_\infty(n)\varepsilon_f(\chi,a,\frac12)$. Clearly
$|\varepsilon_\infty(n)|=1$. We show:
\begin{enumerate}[label=(\alph*)]
\item $\varepsilon_v(\chi,a_v,s)=\varepsilon_v(\chi,a_v,\tfrac12)\,\bigl(\Nm v^{a_v+d_v}\bigr)^{1/2-s}$ for every $v$ and $s$;
\item $\varepsilon_f(\chi,a,s)=\varepsilon_f(\chi,a,\tfrac12)\,\bigl(|d_E|\prod_v\Nm v^{a_v}\bigr)^{1/2-s}$;
\item $|\varepsilon_v(\chi,a_v,\tfrac12)|=1$ for every $v$, hence $|\varepsilon_f(\chi,a,\tfrac12)|=1$.
\end{enumerate}
(a) is the identity $\Nm v^{-s(a_v+d_v)}=\Nm v^{-(a_v+d_v)/2}\,\Nm v^{(a_v+d_v)(1/2-s)}$.
For (b), let $S$ be the finite set of places dividing the different or with $a_v\ne0$. Outside $S$ we have $a_v=d_v=0$ and
$\varepsilon_v=1$ (\lref{GlobalTheory.HeckeEps.locEps_eq_one}), so both products run over $S$, and by (a)
$\varepsilon_f(\chi,a,s)=\varepsilon_f(\chi,a,\frac12)\prod_{v\in S}(\Nm v^{a_v+d_v})^{1/2-s}$. Since the absolute norm of the different
is $|d_E|$, $\prod_{v\in S}\Nm v^{a_v+d_v}=|d_E|\prod_v\Nm v^{a_v}$.
For (c): $|\chi_v|=1$ as $\chi$ is unitary, $|\Nm v^{-(a_v+d_v)/2}\Nm v^{d_v/2}|=\Nm v^{-a_v/2}$, and for $a_v\ge1$ the Gauss sum
attached to an element of valuation $a_v+d_v$ satisfies $|G_v|^2=\Nm v^{a_v}$ because $a_v$ is the conductor exponent
(\lref{GlobalTheory.T.heckeLocalGaussSum_norm_sq_eq_absNorm_pow}; this is the classical evaluation of the absolute value of
a local Gauss sum, \cite[Ch.~2]{Tate1967thesis}). Hence $|\varepsilon_v(\chi,a_v,\frac12)|=1$ for every $v$, and $|\varepsilon_f(\chi,a,\frac12)|=1$.
Inserting (b) in \eqref{eq:pf:sh-div:fe} gives \eqref{eq:pf:sh-div:feabs} with $|\varepsilon|=1$.
\end{proof}

\paragraph{Non-vanishing at integers.}
We use the notation of \cref{sec:pf:fit} for grids: $D$ is a global datum with host $E$ (a CM field with maximal real subfield
$F$), $T$ a set of types with level characters $L_c$ and base type $\mu_0$, and $G$ a grid datum. For $i,j<m$ and a type $t$,
\[
\nu_H(G,t,i,j)=\mu_{0,j}\,\xi_H(t)\,\chi_H(G,m-i)^{-1},
\]
\begin{gather*}
\mathrm{fam}(G,\mu_0,i,j,\mathrm{none})=\nu_H(G,\mu_0,i,j),\\
\mathrm{fam}(G,\mu_0,i,j,c)=\nu_H(G,\mu_0,i,j)\,\delta_c\quad(c\in\{0,1\}),
\end{gather*}
where $\mu_{0,j}$, $\xi_H(t)$ and $\delta_0,\delta_1$ are Hecke characters of $E$ which are part of the level characters and
$\chi_H(G,\cdot)$ the Hecke characters of the grid; $\xi_H(\mu_0)=1$ and every $\xi_H(t)$ has finite order. Let $h(i,j)$ be the half
rank of the square $(i,j)$ (half the sum of the ranks of the spaces at $(i,j)$ and $(i+1,j+1)$) and $e_0(i,j)$ half the
difference of its two bottom ranks; $e_0(i,j)\le h(i,j)$ by a direct comparison of the definitions
(both are defined by explicit formulas in the ranks). The CM type of the host is
$\varphi_D=(\overline{\tau_D})_\tau$, the complex conjugates of the embeddings of $E$ attached to the elements $\tau$ of the CM set
of $D$.

The uniform quantifier $\shdivGBig$ is that of \cref{pf:fit:uniform}: $\shdivGBig(Q)$ asserts that there is $C_{F0}$ such that for
all $\eps$ and $C_F\ge C_{F0}$ there are $C_0,x_0$ such that $Q(G,\mu_0,C,x,b)$ holds for all $C\ge C_0$, all global data $D$ with
size data of exponent $\eps$ and scale $x\ge x_0$ satisfying the standing conditions $\sfname{Rec}$, all level characters with base
type $\mu_0$ and all grid data $G$ at level $(1,C_F,b)$. Two square statements are concerned:
\begin{itemize}
\item $\sfname{FieldDat}=\shdivGBig(Q_1)$, where $Q_1$ asserts: $E$ is a CM field, abelian over $\Q$; the field data
$\sfname{IsFieldData}(d,\varphi_D)$ hold, namely $d$ is a prime $>4$, $E/F$ is unramified at all finite places, $\prod_i\varphi_i(x)=1$
for every root of unity $x\in E$, $2\mid d_E$, $|I|\ge2$, $\varphi_D$ is a CM type, and there is $\omega\in E$ with $\omega^2=-d$ and
$\varphi_i(\omega)=i\sqrt d$ for all $i$; the conductor $\mathfrak f(G)=\prod_{v\in S_0}\mathfrak p_v^{c(v)}$ of the grid ($S_0$ its
set of places of $F$, $c(v_0)=1$ at the base place and $c(v)=2$ otherwise) satisfies the \emph{base condition}: $\mathfrak f\ne0$,
$\mathfrak f\ne\OO_F$, and every prime $v\mid\mathfrak f$ splits in $E$, has exponent $\le2$ in $\mathfrak f$ and lies over an odd
prime $l$ not dividing $d_E$ such that every prime of $E$ above $l$ has norm $l$; and there is a period lattice $\Lambda$ with
$\Lambda=\{x\in\C:\ x$ an algebraic integer of the form $a+bi\sqrt d$, $a,b\in\Q\}$ (an \emph{integral lattice of discriminant $-d$});
\item $\sfname{LNe}=\shdivGBig(Q_2)$, where $Q_2$ asserts: for all $i,j<m$, (a) for every type $t$ there is a Hecke
$L$-function $\Lambda$ of $\nu_H(G,t,i,j)$ with $\Lambda(s)\ne0$ for all integers $s$ with $1-e_0(i,j)\le s\le h(i,j)$, and (b) for
every $c$ there is a Hecke $L$-function $\Lambda$ of $\mathrm{fam}(G,\mu_0,i,j,c)$ with $\Lambda(s)\ne0$ for all integers $s\ge1$.
\end{itemize}

\begin{lemma}[Non-vanishing at $s=1$ and in the half plane of absolute convergence]\label{pf:sh-div:lone}
Let $\chi$ be a unitary Hecke character of a number field $E$ and $\Lambda$ a Hecke $L$-function of $\chi$. Then $\Lambda(s)\ne0$ for
every integer $s\ge1$.
\formallabel{GlobalTheory.CMVal.IsHeckeL.ne_zero_nat}
\end{lemma}

\begin{proof}
For $s\ge2$, $\Lambda(s)$ is the exponential of a convergent series. For $s=1$ this is the classical non-vanishing of
$L(1,\chi)$ \cite{Hecke1920}, \cite[Ch.~VII]{Neukirch1999}; the formal proof
(\lref{GlobalTheory.eulerProduct_ne_zero_one_of_entire}) is the argument with
$\zeta_E(s)^2L(s,\chi)L(s,\bar\chi)$: if $\Lambda(1)=0$, then, with $Z(s)=(s-1)\zeta_E(s)$ entire
(\lref{GlobalTheory.dedekindZeta_eulerProduct_mul_sub_one_entire}), the function
$Z(s)^2\,\frac{\Lambda(s)}{s-1}\,\frac{\overline{\Lambda(\bar s)}}{s-1}$ is entire and equals for $\Re s>1$ the exponential of a
Dirichlet series with non-negative coefficients $\sum_{v,k}(2+2\Re c_v^k)\,k^{-1}\Nm v^{-ks}$; by Landau's theorem on Dirichlet series
with non-negative coefficients this is impossible.
\end{proof}

\begin{lemma}[Non-vanishing at integers left of $1$]\label{pf:sh-div:lnezeroint}
Assume $\sfname{HeckeFunctionalEquationAbs}$ and the conclusion of \cref{pf:sh-div:strips}. Let $E$ be a CM field with a CM type
$\varphi$, and $\chi$ a unitary Hecke character of $E$ of weight $m\in\Z^I$ with $m_i>0$ for all $i$. Then $\chi$ has a Hecke
$L$-function $\Lambda$, and $\Lambda(s)\ne0$ for every integer $s$ with $s+m_i>0$ for all $i$.
\formallabel{GlobalTheory.LValNe.heckeL_ne_zero_int}
\formallabel{GlobalTheory.LValNe.exists_heckeL}
\end{lemma}

\begin{proof}
A Hecke $L$-function of $\chi$ exists because $m\ne0$
(\lref{GlobalTheory.ToricDistribution.exists_isHeckeL_of_hasInf}: a unitary Hecke character of non-zero weight of a CM field has
a Hecke $L$-function). Let $\Lambda$ be any Hecke $L$-function of $\chi$ and $s$ an integer with $s+m_i>0$ for all $i$. If $s\ge1$, use
\cref{pf:sh-div:lone}. Let $s\le0$. The character $\chi$ has infinity type $2m\ne0$ and $\chi^{-1}$ (unitary) has infinity type
$-2m$ and weight $-m\ne0$. By \cref{pf:sh-div:strips} there are completed $L$-functions $\Lambda_c$ of $(\chi,2m)$ and $\Lambda'_c$ of
$(\chi^{-1},-2m)$, and $\chi^{-1}$ has a Hecke $L$-function $\Lambda'$. If $f$ is a Hecke $L$-function and $f_c$ a completed
$L$-function of the same character and infinity type $n$, then $f_c(z)=\Gamma_\infty(n,z)f(z)$ at every $z$ such that no
$z+|n_i|/2$ is a pole of $\Gamma$, by analytic continuation from $\Re z>1$
(\lref{GlobalTheory.LValNe.completed_eq}). At $z=1-s$ we have $z+|{-2m_i}|/2=1-s+m_i\ge1$, so
$\Lambda'_c(1-s)=\Gamma_\infty(-2m,1-s)\Lambda'(1-s)\ne0$ by \cref{pf:sh-div:lone} (as $1-s\ge1$). By \eqref{eq:pf:sh-div:feabs}, in which
$\varepsilon\ne0$ and the base of the power is a non-zero integer, $\Lambda_c(s)=0$ if and only if $\Lambda'_c(1-s)=0$
(\lref{GlobalTheory.Published.HeckeFEModel.zero_iff}); hence $\Lambda_c(s)\ne0$. At $z=s$ we have $s+|2m_i|/2=s+m_i\ge1$, so
$\Lambda_c(s)=\Gamma_\infty(2m,s)\Lambda(s)$, and therefore $\Lambda(s)\ne0$.
\end{proof}

\begin{proposition}[Non-vanishing of the square $L$-functions]\label{pf:sh-div:lne}
$\sfname{HeckeFunctionalEquationAbs}$, the conclusion of \cref{pf:sh-div:strips} and $\sfname{FieldDat}$ imply $\sfname{LNe}$.
\formallabel{GlobalTheory.SqTop.lNe_of_functionalEquation}
\formallabel{GlobalTheory.SqTop.lNe_square}
\formallabel{GlobalTheory.SqTop.squareBottomHalfRankDiff_le_halfRank}
\end{proposition}

\begin{proof}
Take the thresholds $C_{F0}$, and then $C_0,x_0$, of $\sfname{FieldDat}$. In the situation of $\sfname{LNe}$, $\sfname{FieldDat}$ gives
that $E$ is a CM field and that $\varphi_D$ is a CM type; it remains to prove $Q_2$ for fixed $i,j$, and this uses nothing else. Let $\kappa=\kappa(i,j)\in\N^I$ be the weight of the square at $s=0$. By
\lref{GlobalTheory.FunObj.SqNum.Model.isValChar_nu}, $\nu_H(G,\mu_0,i,j)$ and all $\nu_H(G,\mu_0,i,j)\delta_c$ are unitary, trivial on
$\A_F^\times$, of weight $\kappa$ and with prescribed conductor exponents; and by
\lref{GlobalTheory.FunObj.SqNum.Model.kap_add}, $\kappa_\tau(i,j,s)+s$ does not depend on $s\le h(i,j)$ and $\kappa_\tau(i,j,s)>0$,
so $\kappa_\tau=\kappa_\tau(i,j,h(i,j))+h(i,j)>h(i,j)$ for every $\tau$.

(b) For each $c$, $\mathrm{fam}(G,\mu_0,i,j,c)$ is unitary of positive weight $\kappa$, so it has a Hecke $L$-function, which does not
vanish at the integers $s\ge1$ by \cref{pf:sh-div:lone}.

(a) Since $\xi_H(\mu_0)=1$, $\nu_H(G,t,i,j)=\nu_H(G,\mu_0,i,j)\,\xi_H(t)$. The character $\xi_H(t)$ has finite order, hence is trivial
on $E^\times$ at infinity (\lref{GlobalTheory.TypeValues.inf_eq_one}: a Hecke character of finite order of a totally complex field
is trivial on the connected group $\prod_{w\mid\infty}E_w^\times$); thus $\nu_H(G,t,i,j)$ also has weight $\kappa$. It is unitary, being
the character attached to a unitary idele class character of the grid
(\lref{GlobalTheory.FunObj.SqNum.chrEq_nu} and \lref{GlobalTheory.ChrEqUnique.isUnitary_of_chrEq}). By
\cref{pf:sh-div:lnezeroint} it has a Hecke $L$-function $\Lambda$ with $\Lambda(s)\ne0$ whenever $s+\kappa_\tau>0$ for all $\tau$. If
$s\ge1-e_0(i,j)$, then $s+\kappa_\tau>1-e_0(i,j)+h(i,j)\ge1$, since $e_0\le h$.
\end{proof}

\begin{remark}\label{pf:sh-div:rem:lne}
In (a) the upper bound $s\le h(i,j)$ is not used: the proof gives non-vanishing at all integers $s\ge1-e_0(i,j)$. The boundedness
in vertical strips supplied by \cref{pf:sh-div:strips} is not used either; only the existence of completed $L$-functions is.
\end{remark}

\begin{proposition}[Field data]\label{pf:sh-div:fielddat}
$\sfname{FieldDat}$ holds, with all thresholds equal to $0$.
\formallabel{GlobalTheory.SqTop.fieldDat}
\formallabel{GlobalTheory.SqTop.FieldDatPf.isCMType_cmPhi}\formallabel{GlobalTheory.SqTop.FieldDatPf.isBaseCond_condF}
\formallabel{GlobalTheory.SqTop.FieldDatPf.good_prime}\formallabel{GlobalTheory.SqTop.FieldDatPf.splitsC_of_mod}
\end{proposition}

\begin{proof}
Let $D$, size data and a grid $G$ be given as in $\shdivGBig$; we use only that $D$ has a host and that its size data contain
the parameter conditions $g_i\equiv1\bmod Q_Sqd$ ($i=0,1$) for the two level primes and the complete splitting of the primes
$p_{de},p'_{de}$ in $E$. Recall that the host field $E$ is a CM field of degree $4\,\mathrm{ord}$ over $\Q$, Galois over $\Q$ with an
injective homomorphism of $\Gal(E/\Q)$ into an abelian group; hence $E$ is abelian over $\Q$
(\lref{GlobalTheory.ToricDenomBridge.host_isAbelianGalois}).

\emph{The CM type.} By the construction of the host, the embeddings $\varphi_D$ form a CM type and $\varphi_\tau(\sqrt{-d})$ does not
depend on $\tau$, where $\sqrt{-d}\in E$ is the square root of $-d$ of the host
(\lref{GlobalTheory.LamTie.CMType.phi_cmType}, transported along the bijection between the CM set of $D$ and the CM type
of the host).

\emph{Field data.} $d$ is prime and $d\equiv7\bmod16$, so $d>4$; $E/F$ is unramified at the finite places, $2\mid d_E$, and the product
of the $\varphi_i(x)$ over a CM type with constant value at $\sqrt{-d}$ is $1$ for every root of unity $x$
(\lref{GlobalTheory.ToricDenomBridge.host_unramFin}, \lref{GlobalTheory.ToricDenomBridge.two_dvd_discr_host}, and \lref{UpperTheory.prod_CMType_embeddings_root_of_unity_eq_one});
$|I|\ge2$ since $[E:\Q]\le2|I|$ for a CM type and $[E:\Q]=4\,\mathrm{ord}\ge4$. Finally $\varphi_i(\sqrt{-d})^2=-d$, so
$\varphi_i(\sqrt{-d})=\pm i\sqrt d$ with a sign independent of $i$; take $\omega=\pm\sqrt{-d}$ accordingly
(\lref{GlobalTheory.ToricDenomBridge.isFieldData_host}).

\emph{Base condition.} The places of $S_0$ lie over one of the five primes $g_0,g_1,p_{de},p'_{de},q_0$ (the base place lies over
the auxiliary prime $q_0$). Each of them is odd, does not divide $d_E$, and every prime $\mathfrak P$ of $E$ above it has norm $p$ and
lies over a place of $F$ which splits in $E$: for $g_0,g_1$ and $q_0$ this follows from $p\equiv1\bmod Q_Sqd$ (for $q_0$ from its
defining congruence), since $16\mid Q_S$ gives $p\equiv1\bmod4$ and $p$ splits completely in $E$
(by \cref{pf:sh-div:split}, and since a prime which splits completely in $E$ does not divide $d_E$ and the primes above it have norm $p$, \lref{GlobalTheory.SplitEmbSep.not_dvd_discr_of_split} and \lref{GlobalTheory.Mem.Chr.splitsC_main}); for $p_{de},p'_{de}$ from
$p\equiv1\bmod15$ and the complete splitting of the size data. Since $\mathfrak f(G)=\prod_{v\in S_0}\mathfrak p_v^{c(v)}$ with
$c(v)\le2$, every prime dividing $\mathfrak f(G)$ is in $S_0$ and has exponent $\le2$; $\mathfrak f(G)\ne0$, and $\mathfrak f(G)\ne\OO_F$
because the base place has exponent $1$.

\emph{The lattice.} Let $\Lambda=\Z+\Z\frac{1+i\sqrt d}2$. Since $d$ is a prime $\equiv3\bmod4$, $\Lambda$ is the ring of integers of
$\Q(i\sqrt d)$: its elements are integral ($\frac{1+i\sqrt d}2$ is a root of $X^2-X+\frac{d+1}4$), and conversely if $x=a+bi\sqrt d$
($a,b\in\Q$) is integral, then so is $\bar x$, hence $t=2a$ and $n=a^2+db^2$ are integers; $(2db)^2=d(4n-t^2)$ shows that $2db$ is an integer divisible by the prime $d$, so $u=2b\in\Z$;
then $4n=t^2+du^2$ and $d\equiv-1\bmod4$ give $t\equiv u\bmod2$, and $x=\frac{t-u}2+u\frac{1+i\sqrt d}2\in\Lambda$
(\lref{GlobalTheory.ToricDenomBridge.isIntLattice_of}).
\end{proof}

\paragraph{Toric sums and their conjugates.}
Let $E$ be a CM field with maximal real subfield $F$, unramified over $F$ at the finite places, $K\subseteq E$ an imaginary
quadratic field, $d>0$ and $\delta\in K$ with $\delta^2=-d$, $\iota:K\to\C$, and $\varphi$ a CM type of $E$ with
$\varphi_i|_K=\iota$ for all $i$. Let $s\ge1$, $\kappa\in\N^I$ with $\kappa_i\ge1$, $\mathfrak f$ an ideal of $\OO_F$ with the base
condition, $Q$ a finite set of places of $F$ split in $E$ and prime to $\mathfrak f$, $c=\mathfrak f\prod_{q\in Q}\mathfrak p_q$, and
$\chi$ a \emph{valuation character}: unitary, trivial on $\A_F^\times$, of weight $\kappa+s$, with conductor exponent at every prime
$\mathfrak P\mid v$ equal to the exponent of $v$ in $\mathfrak f$. Such a tuple is a \emph{toric datum}. Let
$\shdivPic(c)=\A_{E,f}^\times/E^\times\widehat\OO_c^\times$ be the ring class group of conductor $c$ (finite for $c\ne0$,
\lref{GlobalTheory.CMVal.PicFin.finite_picO}), and for $y\in\A_{E,f}^\times$ put $\lambda(y)=\chi_f(y)|y|^s$, where $\chi_f$ is the
finite part of $\chi$ and $|y|$ the idelic norm. For each class $q$ let $S_q$ be a \emph{lattice sum} of weight $(2s,\kappa)$ of the
lattice of $q$ (an Eisenstein-type sum over the lattice of the class, \cite{LeanSource2026}). A \emph{twisted sum} of $(\chi,\psi)$,
for a character $\psi$ of $\shdivPic(c)$, is a number $T=\sum_q\psi(q)^{-1}\lambda(q)^{-1}S_q(0)$; a \emph{toric sum} is a twisted sum
for $\psi=1$ (the two definitions agree for $\psi=1$). Finally put
$\mathrm{co}(\Omega,T)=C_0\,T/\Omega^{2s|I|+2\sum_i\kappa_i}$ with an explicit constant $C_0=C_0(E,d,2s,\kappa)$.

$\sfname{TwistedTransport}$ is the statement: for all $K,d,\delta,\iota$ as above there is a period $\Omega_0\ne0$ such that
(i) $\mathrm{co}(\Omega_0,T)$ is algebraic for every toric datum, every $\psi$ and every twisted sum $T$ of $(\chi,\psi)$; and
(ii) for every algebraic closure $\Omega$ of $K$ with an embedding $j$ extending $\iota$, every Artin map $\mathrm{art}$ of $K$
with values in $\Gal(\Omega/K)^{\mathrm{ab}}$ satisfying the axioms of global class field theory, every $g\in\Aut(\C)$,
$\sigma\in\Gal(\Omega/K)$ with $j\sigma=gj$ and every finite idele $t$ with $\mathrm{art}(t)^{-1}=\sigma^{\mathrm{ab}}$, there is
$\Omega_1\ne0$ such that for every toric datum, every $(\kappa',\chi')$ conjugate to $(\kappa,\chi)$ under $g$ (that is,
$\kappa'_j=\kappa_i$ whenever $g\varphi_i=\varphi_j$, $\lambda_{\chi'}=g\circ\lambda_\chi$ and $\chi'$ has weight $\kappa'+s$), every $\psi$
and twisted sums $T$ of $(\chi,\psi)$ and $T'$ of $(\chi',g\psi)$:
\[
g\bigl(\mathrm{co}(\Omega_0,T)\bigr)=\mathrm{co}\Bigl(\Omega_1,\ |t|^{-\sum\kappa_i}\bigl(\lambda_{\chi'}(t_E)\,(g\psi)([t_E])\bigr)^{-1}T'\Bigr),
\]
where $t_E$ is the image of $t$ in $\A_{E,f}^\times$ and $[t_E]$ its class.
$\sfname{ToricRatioConjugates}$ is the statement: for every toric datum there are a non-zero toric sum $T_0$ of $\chi$ (for the
conductor $\mathfrak f$) and twisted sums $T(\psi)$ of $(\chi,\psi)$ for all $\psi$, such that every ratio $T(\psi)/T_0$ is algebraic,
and for all $(\Omega,j,\mathrm{art},g,\sigma,t)$ as in (ii) with $g$ fixing all $\varphi_i(x)$ and all values of $\lambda_\chi$, and all
$\psi$,
\begin{equation}\label{eq:pf:sh-div:toricconj}
g\Bigl(\frac{T(\psi)}{T_0}\Bigr)=(g\psi)([t_E])^{-1}\,\frac{T(g\psi)}{T_0},
\end{equation}
and $(g\psi)([t_E])$ is a root of unity.

\begin{proposition}[Conjugates of toric ratios]\label{pf:sh-div:toric}
$\sfname{TwistedTransport}$ implies $\sfname{ToricRatioConjugates}$.
\formallabel{GlobalTheory.ToricTransport.toricRatioConjugates_of_transport}
\formallabel{GlobalTheory.ToricTransport.exists_toricSum_ne_zero}\formallabel{GlobalTheory.ToricTransport.exists_twistedSum}
\formallabel{GlobalTheory.ToricTransport.isTwistedSum_one}
\end{proposition}

\begin{proof}
Fix a toric datum, and let $\Omega_0$ be the period of $\sfname{TwistedTransport}$ for $(K,d,\delta,\iota)$.

\emph{Existence.} $\chi$ has a Hecke $L$-function $\Lambda$ (\cref{pf:sh-div:lnezeroint}), and the toric sum attached to
the empty set of auxiliary places is the value $\Lambda(s)$ (\lref{GlobalTheory.ToricSumValue.toricSum_eq}), which is non-zero by
\cref{pf:sh-div:lone} since $s\ge1$; let $T_0=\Lambda(s)$. For every
$\psi$ a twisted sum $T(\psi)$ exists: $\shdivPic(c)$ is finite, the finite part of $\chi$ is trivial on the units of the order
of conductor $c$, and every class has a lattice sum
(\lref{GlobalTheory.CMVal.OrderUnits.finPart_eq_one} and \lref{GlobalTheory.CMVal.LatEnt.exists_isLatSum}).

\emph{Algebraicity.} The normalising factors of $\mathrm{co}(\Omega_0,\cdot)$ do not depend on $\psi$, so
$T(\psi)/T_0=\mathrm{co}(\Omega_0,T(\psi))\,\mathrm{co}(\Omega_0,T_0)^{-1}$ (\lref{GlobalTheory.ToricTransport.coordIn_div}),
a quotient of algebraic numbers by (i), applied to $(\chi,\psi)$ and to the toric datum with $Q=\emptyset$ and $\psi=1$.

\emph{Conjugation.} Let $g$ fix all $\varphi_i(x)$ and all values of $\lambda_\chi$, and let $\Omega_1$ be given by (ii). Then
$(\kappa,\chi)$ is conjugate to itself under $g$: if $g\varphi_i=\varphi_j$ then $\varphi_i=\varphi_j$ (as $g$ fixes $\varphi_i(E)$),
hence $i=j$ by the uniqueness in the definition of a CM type; and $\lambda_\chi=g\circ\lambda_\chi$. Apply (ii) to $T=T(\psi)$,
$T'=T(g\psi)$, and to $T=T'=T_0$ with $\psi=1$ (note $g\cdot1=1$). The common factors $|t|^{-\sum\kappa_i}$ and $\lambda_\chi(t_E)^{-1}$
are non-zero (\lref{GlobalTheory.IndefiniteRowSideFullHecke.LatticeNormTransport.lamF_ne_zero}) and cancel in the quotient, which
gives \eqref{eq:pf:sh-div:toricconj}. Finally $(g\psi)([t_E])$ is a value of a character of the finite group $\shdivPic(c)$, hence a
root of unity.
\end{proof}

%==========================================================================================================================
\subsection{Primes in arithmetic progressions and split primes}\label{sec:pf:sh-div:primes}

For $M\ge1$ let $\pi(X;M,b)$ be the number of primes $p\le X$ with $p\equiv b\bmod M$, $\psi(X;M,b)=\sum_{n\le X,\,n\equiv b\,(M)}\Lambda(n)$,
$\shdivLi(X)=\int_2^X\frac{dt}{\log t}$, and for $\theta\in\Z$, $\beta\in\R$
\[
\shdivEm(\theta,\beta,X)=\shdivLi(X)-\theta\,\shdivLi(X^\beta),\qquad
\rho(C,c,B,X)=Ce^{-c\log X/\log B}+Ce^{-c\sqrt{\log X}} .
\]
For $c_1>0$ and $M\ge3$, a pair $(\chi,\beta)$ of a Dirichlet character modulo $M$ and a real number is \emph{exceptional} if
$\chi^2=1$, $1-c_1/\log M<\beta<1$ and $L(\beta,\chi^*)=0$, where $\chi^*$ is the primitive character inducing $\chi$. For $b$ prime
to $M$, $(\theta,\beta)$ is an \emph{exceptional datum} for $b$ if either $\theta=0$ and there is no exceptional pair, or there is
an exceptional pair $(\chi,\beta)$ with $\chi(b)=\theta$.

$\sfname{PrimesRel}$ is the statement: there are $c_1\in(0,\frac18]$ and $c_2,c_3,C_A,c_S>0$ such that for every $M\ge3$:
\begin{enumerate}[label=(P\arabic*)]
\item any two exceptional pairs modulo $M$ are equal;
\item every exceptional pair satisfies $1-\beta\ge c_S/(M^2\log M)$;
\item for $b$ prime to $M$, every exceptional datum $(\theta,\beta)$ for $b$ and every $X\ge M^{c_2}$,
\[
\bigl|\varphi(M)\pi(X;M,b)-\shdivEm(\theta,\beta,X)\bigr|\le C_A\bigl(e^{-c_3\log X/\log M}+e^{-\sqrt{c_3\log X}}\bigr)\shdivEm(\theta,\beta,X).
\]
\end{enumerate}

\begin{theorem}[Primes in progressions with relative error]\label{pf:sh-div:primesrel}
$\sfname{PrimesRel}$ holds.
\formallabel{GlobalTheory.Lin.primesRel_holds}
\formallabel{GlobalTheory.Lin.sign_of_sq_eq_one}
\end{theorem}

(P3) is a form of the prime number theorem for arithmetic progressions in the range $X\ge M^{c_2}$, with the contribution of a
possible exceptional zero; estimates of this kind go back to Gallagher \cite{Gallagher1970} (see also \cite[Ch.~18]{IwaniecKowalski2004},
where they are used for Linnik's theorem). (P1) is the Landau--Page theorem \cite[Ch.~14]{Davenport2000}, and (P2) is much weaker
than the classical bound $1-\beta\gg M^{-1/2}(\log M)^{-2}$ of loc.\ cit. In the formal statement the error term is relative to the
main term, uniformly in $X\ge M^{c_2}$, and primes are counted directly (rather than weighted by $\Lambda$).

\begin{proof}
The formal proof combines three results, each a formalized classical theorem:
\begin{enumerate}[label=(\alph*)]
\item (\lref{GlobalTheory.T.exceptionalDirichletZero_unique_zeroFreeRegion}) There is $a>0$ such that for every
$c_1\in(0,a]$ there is $c>0$ such that for all $M\ge3$: exceptional pairs are unique; an exceptional pair $(\chi,\beta)$ has $\chi\ne1$,
$\beta\ge\frac78$, $L(\beta,\chi)=0$ with a simple zero, and $1-\beta\ge c/(M^2\log M)$; and every zero $\rho$ of $L(s,\chi)$ with
$\Re\rho<1$ is either real with $(\chi,\rho)$ exceptional or satisfies $\Re\rho\le1-c/\log(M(|\Im\rho|+2))$. The formal proof uses
a lower bound for $|L(1,\chi)|$ (\lref{GlobalTheory.T.exists_dirichletLFunction_zero_gap_near_one}).
\item (\lref{GlobalTheory.T.chebyshevPsiProgression_relative_error}) There is $a'>0$ such that for every $c_1\in(0,a']$ there
are $c_2,c_3,C>0$ with
$\bigl|\varphi(M)\psi(u;M,b)-(u-\theta u^\beta/\beta)\bigr|\le C\bigl(e^{-c_3\log u/\log M}+e^{-\sqrt{c_3\log u}}\bigr)(u-\theta u^\beta/\beta)$
for all $M\ge3$, $b$ prime to $M$, exceptional data $(\theta,\beta)$ for $b$ and $u\ge M^{c_2}$. This is Gallagher's theorem; the formal
proof follows his method: the explicit formula, a log-free zero density estimate, the Deuring--Heilbronn repulsion of the zeros
by an exceptional zero, and bounds for higher logarithmic derivatives
(\lref{GlobalTheory.T.logFree_zeroDensity_dirichlet_one_modulus}, \lref{GlobalTheory.T.explicit_formula_dirichlet_zeros_right_half}, and \lref{GlobalTheory.T.one_sub_re_ge_of_exceptional_zero_repulsion}).
\item (\lref{GlobalTheory.T.primeCount_relative_error_of_psi_relative_error}) Partial summation: for all $c_2,c_3,C,c_S>0$
there are $c'_2,c'_3,C'>0$ such that if $B\ge3$, $1\le\phi\le B$, $\theta=0$ or ($\theta=\pm1$, $\beta\ge\frac78$ and $c_S/B^3\le1-\beta$),
and a set $P$ of integers satisfies the estimate of (b) (with $\phi$, $B$ in place of $\varphi(M)$, $M$) for $\sum_{n\le u,\,n\in P}\Lambda(n)$,
then $|\phi\,\#\{p\le X\ \text{prime},\ p\in P\}-\shdivEm(\theta,\beta,X)|\le C'(e^{-c'_3\log X/\log B}+e^{-\sqrt{c'_3\log X}})\shdivEm(\theta,\beta,X)$
for $X\ge B^{c'_2}$.
\end{enumerate}
Let $c_1=\min(\frac18,a,a')$, let $c$ be the constant of (a) and $c_2,c_3,C$ those of (b) for $c_1$, and $c'_2,c'_3,C'$ those of (c) for
$(c_2,c_3,C,c)$. We claim that $\sfname{PrimesRel}$ holds with $(c_1,c'_2,c'_3,C',c_S=c)$. (P1) and (P2) are part of (a). For (P3), let
$b$ be prime to $M$ and $(\theta,\beta)$ an exceptional datum. If $\theta\ne0$ there is an exceptional pair $(\chi,\beta)$ with
$\chi(b)=\theta$; since $\chi^2=1$ and $b$ is a unit, $\theta=\chi(b)=\pm1$, $\beta\ge\frac78$ by
(a), and $c/M^3\le c/(M^2\log M)\le1-\beta$ as $\log M\le M$. Apply (c) with $B=M$, $\phi=\varphi(M)\in[1,M]$ and $P$ the residue
class of $b$, whose hypothesis is (b).
\end{proof}

$\sfname{PrimesAPInterval}(C_0,C_r,c_r,c)$ is the statement: for all $m\ge3$, $a$ prime to $m$ and reals $Y,h$ with
$Y\ge(4m)^{C_0}$, $h\le1$, $\rho(C_r,c_r,4m,Y)\le h/100$ and $(4m)^7\le hY^{1/3}$,
\begin{equation}\label{eq:pf:sh-div:apinterval}
\#\{p\ \text{prime}:\ Y<p\le Y(1+h),\ p\equiv a\bmod m\}\ \ge\ c\,h\,(4m)^{-6}\,\frac{Y}{\log Y}.
\end{equation}

\begin{lemma}[Short intervals from relative asymptotics]\label{pf:sh-div:shortint}
For all $C_0,C_r,c_r,c_b>0$ there are $C'_0\ge C_0$ and $c>0$ with the following property. Let $B\ge4$, $1\le G\le B$,
$\theta\in\{-1,0,1\}$ and $\beta\in\R$ with $\frac12\le\beta<1$ and $c_b/(B^2\log B)\le1-\beta$ if $\theta\ne0$; let $\pi_1,\pi_2$ be
$\N$-valued functions with $\pi_2(Y)\le2\sqrt Y$ for $Y\ge1$, and $\eta$ a function with $0\le\eta\le3\log B$, such that
\begin{equation}\label{eq:pf:sh-div:relasym}
\bigl|G(\pi_1(Y)+\pi_2(Y)+\eta(Y))-\shdivEm(\theta,\beta,Y)\bigr|\le\rho(C_r,c_r,B,Y)\,\shdivEm(\theta,\beta,Y)\qquad(Y\ge B^{C_0}).
\end{equation}
Then for $X\ge B^{C'_0}$: $\shdivEm(\theta,\beta,X)\ge cB^{-5}X/\log X$, and for every $h\le1$ with $\rho(C_r,c_r,B,X)\le h/100$ and
$B^7\le hX^{1/3}$, $\pi_1(X(1+h))-\pi_1(X)\ge c\,h\,B^{-6}X/\log X$.
\formallabel{Formalism.HermitianLattices.GlobalBoundRefinements.shortIntervalLowerBound}
\formallabel{Formalism.HermitianLattices.GlobalBoundRefinements.exceptional_zero_contribution_bound}
\formallabel{Formalism.HermitianLattices.GL2PrimeSums.Emain_lower_bound}
\formallabel{Formalism.HermitianLattices.GlobalBoundRefinements.exists_adjusted_exceptional_exponent}
\formallabel{Formalism.HermitianLattices.GL2PrimeSums.gl2a4_weight}
\end{lemma}

\begin{proof}
Take $C'_0=\max(C_0,60,2/c_b)$ and $c=\frac1{100}$, and let $X\ge B^{C'_0}$. Write $u=X^{1/12}$; then $u\ge1024$ (as $X\ge4^{60}$),
$\sqrt X=u^6$, $X^{1/3}=u^4$ and $\log X\le12u$; put $L=\log X$. Note $\frac{d}{dt}\shdivEm(\theta,\beta,t)=(1-\theta t^{\beta-1})/\log t$.

\emph{Normalising $\beta$.} If $\theta=0$, $\shdivEm$ does not depend on $\beta$ and we replace $\beta$ by $\frac12$. If $\theta\ne0$, then
$L\ge\frac2{c_b}\log B$ gives $(1-\beta)L\ge c_bL/(B^2\log B)\ge2/B^2$. So in all cases
$\frac12\le\beta<1$, and $(1-\beta)L\ge2/B^2$ if $\theta=1$.

\emph{Weights.} For $Z\ge1$ such that $(1-\beta)\log Z\ge1/B^2$ when $\theta=1$, there is $W$ with $\frac1{2B^2}\le W\le1$,
$W\le1-\theta t^{\beta-1}$ for $t\ge Z$ and $1-\theta t^{\beta-1}\le2W$ for $1\le t\le Z^2$
(take $W=1$ if $\theta\ne1$; if $\theta=1$ take $W=1-Z^{\beta-1}$: then $W\le1-t^{\beta-1}$ for $t\ge Z$, $1-t^{\beta-1}\le1-Z^{2(\beta-1)}=W(1+Z^{\beta-1})\le2W$ for $t\le Z^2$, and $W\ge x/(1+x)\ge1/(2B^2)$ for $x=(1-\beta)\log Z\ge1/B^2$). Integrating
the derivative, $W(b-a)/\log b\le\shdivEm(b)-\shdivEm(a)$ if $W\le1-\theta t^{\beta-1}$ on $[a,b]$, and
$\shdivEm(b)-\shdivEm(a)\le V(\shdivLi(b)-\shdivLi(a))$ if $1-\theta t^{\beta-1}\le V$ on $[a,b]$; moreover $0\le\shdivEm(Y)$ for $Y\ge4$ and
$\shdivEm(4)\le6$, and $\shdivLi(Y)\le3Y/\log Y$ for $Y\ge1296$.

\emph{Main term.} With $Z=\sqrt X$ (so $(1-\beta)\log Z\ge1/B^2$ if $\theta=1$): $\shdivEm(X)\ge\shdivEm(X)-\shdivEm(\sqrt X)\ge W(X-\sqrt X)/L\ge
WX/(2L)\ge X/(4B^2L)\ge\frac1{100}B^{-5}X/L$.

\emph{Short intervals.} Let $h$ be as in the statement; $h>0$ since $\rho>0$. Put $X'=X(1+h)\le2X\le X^2$, so $\log X'\le2L$, and take
the weight $W$ for $Z=X$; let $T=WhX/L$ and $\Delta=\pi_1(X')-\pi_1(X)$. Then:
(i) $\shdivEm(X')-\shdivEm(X)\ge WhX/\log X'\ge T/2$;
(ii) $\shdivEm(X')\le\shdivEm(4)+2W(\shdivLi(X')-\shdivLi(4))\le6+12WX/L$, so $h\,\shdivEm(X')\le6h+12T$;
(iii) by \eqref{eq:pf:sh-div:relasym} at $X'$ and at $X$, and since $\rho(\cdot,\cdot,B,Y)$ decreases in $Y$, the value of
$G(\pi_1+\pi_2+\eta)$ is $\ge(1-\frac h{100})\shdivEm(X')$ at $X'$ and $\le\shdivEm(X)+\frac h{100}\shdivEm(X')$ at $X$;
(iv) $G\pi_2(X')\le B\cdot2\sqrt{X'}\le3Bu^6$ and $G\eta(X')\le3B\log B\le Bu$;
(v) from $B^7\le hu^4$ and $L\le12u$: $4Bu^6\le\frac{12}{100}\frac{hX}{2B^2L}\le\frac{12}{100}T$ and $\frac{12}{100}h+Bu\le Bu^6$.
Writing $G\Delta$ as the difference of the values of $G(\pi_1+\pi_2+\eta)$ at $X'$ and $X$ minus $G(\pi_2+\eta)(X')$ plus
$G(\pi_2+\eta)(X)\ge0$, (i)--(v) give
\[
G\Delta\ge\frac T2-\frac1{50}(6h+12T)-3Bu^6-Bu\ge\frac T2-\frac{24}{100}T-4Bu^6\ge\frac{14}{100}T .
\]
Hence $\Delta\ge0$, and as $G\le B$ and $W\ge1/(2B^2)$, $\Delta\ge\frac{14}{100}\frac T B\ge\frac7{100}\frac{hX}{B^3L}\ge\frac1{100}hB^{-6}\frac XL$.
\end{proof}

\begin{theorem}[Primes in progressions in short intervals]\label{pf:sh-div:apinterval}
$\sfname{PrimesRel}$ implies that there are $C_0,C_r,c_r,c>0$ such that $\sfname{PrimesAPInterval}(C_0,C_r,c_r,c)$ holds.
\formallabel{GlobalTheory.Lin.primesAPInterval_of_primesRel}
\formallabel{Formalism.HermitianLattices.GlobalBoundRefinements.shortIntervalImpliesProgressions_statement}
\end{theorem}

\begin{proof}
Let $(c_1,c_2,c_3,C_A,c_S)$ be the constants of $\sfname{PrimesRel}$ and $c_r=\min(c_3,\sqrt{c_3})>0$.

\emph{Step 1: a relative asymptotic with modulus $4m$.} Let $m\ge3$ and $a$ prime to $m$. If there is an exceptional pair
$(\chi,\beta)$ modulo $m$, put $\theta=\chi(a)$; then $\theta=\pm1$ (as $\chi^2=1$), $\beta\ge1-c_1/\log m\ge\frac12$ (since
$\log m>1$ and $c_1\le\frac18$), $\beta<1$, and by (P2)
$c_S/((4m)^2\log4m)\le c_S/(m^2\log m)\le1-\beta$. Otherwise put $\theta=0$, $\beta=\frac12$. In both cases $(\theta,\beta)$ is an
exceptional datum for $a$. For $Y\ge(4m)^{c_2}\ge m^{c_2}$, (P3) gives
$|\varphi(m)\pi(Y;m,a)-\shdivEm(\theta,\beta,Y)|\le\rho(C_A,c_r,4m,Y)\shdivEm(\theta,\beta,Y)$, because
$e^{-c_3\log Y/\log m}\le e^{-c_r\log Y/\log4m}$ and $e^{-\sqrt{c_3\log Y}}\le e^{-c_r\sqrt{\log Y}}$ (and the inequality of (P3)
forces $\shdivEm(\theta,\beta,Y)\ge0$). This is the hypothesis
\eqref{eq:pf:sh-div:relasym} of \cref{pf:sh-div:shortint} for $B=4m$, $G=\varphi(m)$ (note $1\le\varphi(m)\le m\le4m$),
$\pi_1=\pi(\cdot\,;m,a)$, $\pi_2=0$, $\eta=0$, $C_0=c_2$, $(C_r,c_r)=(C_A,c_r)$ and $c_b=c_S$.

\emph{Step 2.} Let $C'_0,c$ be given by \cref{pf:sh-div:shortint} for $(c_2,C_A,c_r,c_S)$. For $m,a,Y,h$ as in
$\sfname{PrimesAPInterval}$, the lemma (applied with $X=Y$, $B=4m$) gives \eqref{eq:pf:sh-div:apinterval}. Hence
$\sfname{PrimesAPInterval}(C'_0,C_A,c_r,c)$ holds.
\end{proof}

\paragraph{Split primes in the host.} Recall (\cref{sec:hosts}) that a host $H$ has a CM field $E$, Galois and abelian over $\Q$,
containing the imaginary quadratic field $K=\Q(\sqrt{-d})$, with invariants $Q_S$ (a power of $2$ divisible by $16$), primes $q$ and $d$
($d\equiv7\bmod16$), the conductor $N_H=Q_Sqd$, and a class $\mathrm{cls}(p)$ attached to every prime $p$: $\mathrm{ram2}$ if $p\mid Q_S$,
$\mathrm{ramQ}$ if $p=q$, $\mathrm{ramD}$ if $p=d$, and otherwise $\mathrm{splitK}$ if the Jacobi symbol $(\frac{-d}{p})$ is $1$
(followed by two further classes which do not matter here). The host also records the reciprocity law of $E$: for $p\nmid N_H$ and a
prime $\mathfrak P$ of $E$ above $p$, the inertia group at $\mathfrak P$ is trivial and the arithmetic Frobenius at $\mathfrak P$ is
the automorphism attached to the residue class of $p$ modulo $N_H$.

\begin{proposition}[Primes $\equiv1$ modulo the conductor split]\label{pf:sh-div:split}
Let $H$ be a host and $p$ a prime with $p\equiv1\bmod Q_Sqd$. Then $p$ splits completely in $E$ (the number of places of $E$
above $p$ is $[E:\Q]$) and $\mathrm{cls}(p)=\mathrm{splitK}$.
\formallabel{GlobalTheory.HostSplit.split_of_one_mod}
\formallabel{GlobalTheory.HostSplit.ncard_primesOver_of_stabilizer}
\end{proposition}

\begin{proof}
For every divisor $m>1$ of $Q_Sqd$ we have $p\equiv1\bmod m$; in particular $p\equiv1\bmod4$ (as $4\mid Q_S$) and $p\equiv1\bmod d$, and
$p$ divides no divisor of $Q_Sqd$ (otherwise $p\equiv1\bmod p$). Hence $p\nmid N_H$, $p\nmid Q_S$, $p\ne q$ and $p\ne d$.

\emph{Splitting.} Let $\mathfrak P$ be a prime of $\OO_E$ above $p$. By the reciprocity law of the host, the arithmetic Frobenius at
$\mathfrak P$ is the automorphism attached to $p\equiv1\bmod N_H$, i.e.\ the identity, and the inertia group at $\mathfrak P$ is trivial.
Every element of the decomposition group of $\mathfrak P$ is the product of an element of the inertia group and a power of a Frobenius
element (\lref{Formalism.FrobeniusOrbits.Decomp.stab_eq_inertia_mul}), so the decomposition group is trivial, and the number
of primes above $p$, which is the index of the decomposition group in $\Gal(E/\Q)$, is $[E:\Q]$. Places of $E$ above $p$ correspond to primes of $\OO_E$ above $p$
(\lref{GlobalTheory.HostSplit.ncard_placesOver}).

\emph{The class.} Since $d$ is odd, $p\equiv1\bmod4$ and $p\equiv1\bmod d$, quadratic reciprocity gives
$\bigl(\frac{-d}{p}\bigr)=\bigl(\frac{-1}{p}\bigr)\bigl(\frac dp\bigr)=\bigl(\frac pd\bigr)=\bigl(\frac1d\bigr)=1$
(\lref{GlobalTheory.HostSplit.jacobiSym_neg_eq_one}, for the Jacobi symbol). As $p\nmid Q_S$, $p\ne q$ and $p\ne d$, the
definition of the class gives $\mathrm{cls}(p)=\mathrm{splitK}$.
\end{proof}

%==========================================================================================================================
\subsection{Volumes of unitary groups}\label{sec:pf:sh-div:vol}

Let $E/F$ be a CM extension of number fields with complex conjugation $\sigma$, $M\ge1$ and $J\in M_M(E)$ a non-degenerate hermitian
matrix. Let $H=\Uni(J)$ be its unitary group over $F$, $H(\A_F)$ its group of adelic points (a locally compact second countable
group) and $H(F)\subseteq H(\A_F)$ the rational points (a countable discrete subgroup).

\begin{theorem}[Finite volume of the adelic unitary quotient]\label{pf:sh-div:finvol}
For every Haar measure $\mu$ on $H(\A_F)$ and every measurable fundamental domain $\mathcal F$ for the action of $H(F)$ on
$H(\A_F)$, $\mu(\mathcal F)<\infty$.
\formallabel{GlobalTheory.Mem.VolFin.finVol}
\formallabel{GlobalTheory.Mem.VolFin.stabHat_conj}
\end{theorem}

This is a special case of the theorem of Borel and Harish-Chandra \cite{BorelHarishChandra1962} on the finiteness of the volume of
$G(F)\backslash G(\A_F)$ for reductive groups $G$ without non-trivial $F$-rational characters (the abelianisation of $H$ is the
norm-one torus of $E/F$, which is anisotropic over $F$); see also \cite{Borel1963,PlatonovRapinchuk1994}. The formal statement is
exactly this special case. The formal proof is the classical reduction to the finiteness of class numbers and to arithmetic
subgroups of the real points.

\begin{proof}
Write $H(\A_F)=H_\infty\times H_f$, as topological groups, compatibly with the diagonal embedding of $H(F)$
(\lref{GlobalTheory.Mem.VolSplit.exists_adelic_split}). Let $L=\OO_E^M$ and let $K\subseteq H_f$ be the stabiliser of the
closure $\widehat L$ of $L$; it is compact and open (\lref{GlobalTheory.T.completedLattice_stabilizer_compact_open}).
\begin{enumerate}[label=(\roman*)]
\item \emph{Finiteness of the class number}: there is a finite set $S\subseteq H_f$ with $H_f=\bigcup_{s\in S}H(F)sK$
(\lref{GlobalTheory.T.finite_classSet_hermitianUnitary_openLevel}; \cite{Borel1963}).
\item For $s\in S$ let $\Lambda_s\subseteq E^M$ be the lattice with $\widehat\Lambda_s=s\widehat L$
(\lref{GlobalTheory.AdLat.exists_lattice_latOf}) and $\Gamma_s\subseteq H(F)$ its stabiliser. \emph{Finite covolume of
arithmetic subgroups of the real points}: there is a measurable $\Omega_s\subseteq H_\infty$ of finite Haar measure with
$H_\infty=\Gamma_s\Omega_s$ (\lref{GlobalTheory.T.finite_covolume_latticeStabiliser_unitaryRealPoints}; \cite{BorelHarishChandra1962};
the formal proof constructs, from a CM type and a diagonalisation of $J$ at the real places, an explicit set of finite measure
covering a fundamental domain).
\end{enumerate}
Let $g=(g_\infty,g_f)\in H(\A_F)$. By (i), $g_f=\gamma sk$ with $\gamma\in H(F)$, $s\in S$, $k\in K$. By (ii),
$\gamma^{-1}g_\infty=\delta w$ with $\delta\in\Gamma_s$ and $w\in\Omega_s$. Since $\delta$ stabilises $\Lambda_s$, $s^{-1}\delta s$
stabilises $\widehat L$, i.e.\ lies in $K$. For $\eta=\delta^{-1}\gamma^{-1}\in H(F)$ we
get $(\eta g)_\infty=w\in\Omega_s$ and $(\eta g)_f=\delta^{-1}sk=s(s^{-1}\delta^{-1}s)k\in sK$. Hence $H(\A_F)=H(F)\cdot\bigcup_{s\in S}\Omega_s\times
sK$, where the union has finite Haar measure ($K$ is compact). Since $H(F)$ is countable and $\mu$ is left invariant, every measurable
fundamental domain has measure at most that of the union (\lref{GlobalTheory.Mem.VolFin.measure_lt_top_of_cover}), hence
finite.
\end{proof}

\paragraph{The volume ratio.} We use the notation of \cref{sec:pf:stageU}. For a global datum $D$ with host $E$ (with maximal real
subfield $F$, $f=[F:\Q]$) the development attaches to every hermitian space $W$ over $E$ with a standard lattice a volume $\vol_D(W)$;
$\sfname{VolTie}(D)$ asserts that for every such $W$ of constant signature, every Haar measure in the normalisation of $D$ and every
fundamental domain $\mathcal F$ of $\Uni(W)(F)$ in $\Uni(W)(\A_F)$, the measure of $\mathcal F$ is $\vol_D(W)$. The \emph{volume ratio}
$\sfname{VolumeRatio}_2(D)$ is the statement (VR) of \cref{sec:pf:stageU}: for every hermitian space $W$ of rank $n\ge3$ and signature
$(P,Q)$ at every complex embedding, and every sign $\pm$ with $P\ge1$ (for $+$) or $Q\ge1$ (for $-$), there are $x\in W$ with
$\langle x,x\rangle=\pm1$ and a hermitian space $W'$ of signature $(P-1,Q)$ (resp.\ $(P,Q-1)$) with an isometric embedding onto
$x^\perp$ such that
\begin{equation}\label{eq:pf:sh-div:vr}
\frac{\vol_D(W')}{\vol_D(W)}=\Bigl(\frac{(2\pi)^n}{(n-1)!}\Bigr)^{f}(2v)^{f(\frac12-n)}\,L(n,\eta^n)^{-1},
\end{equation}
with $v$ and $L(n,\eta^n)$ as in \cref{prop:families:mass}. $\sfname{MassRatioUnderVolume}(D)$ is the implication
$\sfname{VolTie}(D)\Rightarrow\sfname{VolumeRatio}_2(D)$. Finally $\sfname{SW1b2}(F,E,\sigma,M,J)$ is the one-vector mean value formula of
\cref{sec:pf:stageU} (item 2 of the proof of \cref{pf:stageU:volratio}): for $M\ge3$, $\beta\in F^\times$, measurable
$\Phi_\infty\ge0$ on $(E\otimes\R)^M$ and non-negative pure finite test functions $\Phi_f=\bigotimes_v\Phi_v$, with self-dual measures
and fibre measures $\mu_{\infty,t},\mu_{v,t}$ of the hermitian norm, the local integrals $\int\Phi_v\,d\mu_{v,\beta}$ are multipliable and
\[
\int_{H(F)\backslash H(\A_F)}\sum_{x\in E^M\smallsetminus0,\ \langle x,x\rangle=\beta}\Phi(h^{-1}x)\,dh
=\int\Phi_\infty\,d\mu_{\infty,\beta}\cdot\prod_v\int\Phi_v\,d\mu_{v,\beta}
\]
(the left side being the mean over a fundamental domain).

\begin{theorem}[The volume ratio from the one-vector formula]\label{pf:sh-div:volratio}
Assume the one-vector formula $\sfname{SW1b2}(F,E,\sigma,M,J)$ for every CM extension $E/F$ and every non-degenerate hermitian $J$.
Then for every global datum $D$, $\sfname{VolTie}(D)$ implies $\sfname{VolumeRatio}_2(D)$; that is, the statement
$\sfname{MassRatioUnderVolume}(D)$ holds.
\formallabel{GlobalTheory.VolRatioFromSecondClause.volumeRatio}
\formallabel{GlobalTheory.VolRatioUnfold.volumeRatio_of_unfolding}
\end{theorem}

\begin{proof}
The proof uses six statements of the development, each proved there unconditionally; we state them in the form used.
\begin{enumerate}[label=(V\arabic*)]
\item\label{pf:sh-div:V:perp} \emph{Unit vectors with self-dual complements}
(\lref{GlobalTheory.UnitPerpObj.unitVectorPerp}): for $W$, $(P,Q)$ and the sign as in (VR) there are $x\in W$ with
$\langle x,x\rangle=\pm1$, a self-dual lattice $L\ni x$ of $W$, and $W'$ of the complementary signature with an isometric
embedding $\iota$ onto $x^\perp$ such that $\iota$ maps the standard lattice of $W'$ onto $L\cap x^\perp$.
\item\label{pf:sh-div:V:indep} \emph{Independence of the lattice} (\lref{GlobalTheory.VolIndepClosed.volIndep}): the volume
of $W$ does not change if its standard lattice is replaced by another self-dual lattice, and is invariant under isometries.
\item\label{pf:sh-div:V:dens} \emph{Local densities of self-dual lattices}
(\lref{GlobalTheory.LatticeDensityClosed.selfDualLatticeDensity}): at a finite place $v$ of $F$, for a self-dual lattice $L_v$
of a non-degenerate hermitian space of rank $M$ over $E_v$ and self-dual measures, $\mathrm{vol}(L_v)=\mathfrak d_v^{-M}$,
$\mathrm{vol}(\OO_{F_v})=\mathfrak d_v^{-1/2}$ ($\mathfrak d_v$ the local dual index), and for $|\beta|_v=1$ the fibre measure of
$L_v$ at $\beta$ is an explicit value $\delta(\mathfrak d_v,q_v,\eta_v,M)$ depending on the splitting type $\eta_v$ of $v$ in $E$, locally
constant in $\beta$.
\item\label{pf:sh-div:V:euler} \emph{The Euler product of the densities}
(\lref{GlobalTheory.DensityProductClosed.eulerProductOfDensities}): for $E/F$ unramified at the finite places, the product over
$v$ of these densities, for $M\ge2$, converges and equals an explicit expression in $|d_F|$ and $L(M,\eta^M)^{-1}$.
\item\label{pf:sh-div:V:real} \emph{The real places} (\lref{GlobalTheory.RealPlaceMeasuresProof.realPlaceMeasures}): the
self-dual measures at the real places are explicit, and for the signature $(P,Q)$ there is a non-negative measurable archimedean test
function $\Phi_\infty$ whose fibre integral at $\pm1$ and whose orbital integral over the quotient of the real unitary group by the
stabiliser of a vector of norm $\pm1$ (computed through the compact dual) are explicit
(\lref{GlobalTheory.VolRatioUnfold.arch}).
\item\label{pf:sh-div:V:unf} \emph{Unfolding} (\lref{GlobalTheory.UnfoldingObj.unfolding}): for $x$ of norm $\beta$ in the
setting of \ref{pf:sh-div:V:perp}, the mean over $H(F)\backslash H(\A_F)$ of $h\mapsto\sum_{\langle y,y\rangle=\beta}\Phi(h^{-1}y)$, for
$\Phi=\Phi_\infty\otimes\mathbf 1_{\widehat L}$, is expressed through the volumes of $W$ and $W'$ and the orbital integral of $\Phi_\infty$,
with the quotient measures of unimodular groups; its formal proof uses that the stabiliser of a self-dual lattice acts transitively
on the lattice vectors of a given unit length (\lref{GlobalTheory.T.stabSelfDualLattice_transitive_on_length}).
\end{enumerate}
Let $D$ satisfy $\sfname{VolTie}(D)$, and let $W,(P,Q)$ and the sign be given. Take $x,L,W',\iota$ from \ref{pf:sh-div:V:perp}; these
are the $x$, $W'$ and embedding required by (VR). Let $W_1$ be $W$ with the lattice $L$; by \ref{pf:sh-div:V:indep},
$\vol_D(W)=\vol_D(W_1)$, and $L\cap x^\perp=\iota(L_{W'})$. It remains to prove \eqref{eq:pf:sh-div:vr} for $W_1$ and $W'$
(\lref{GlobalTheory.VolRatioUnfold.ratio_of_setting}). Apply $\sfname{SW1b2}$ to $W_1$, $\beta=\pm1$, the pure finite test function
$\Phi_f=\mathbf 1_{\widehat L}=\bigotimes_v\mathbf 1_{L_v}$ and the archimedean function $\Phi_\infty$ of \ref{pf:sh-div:V:real}. On the right
side, the local integrals at the finite places are the densities of \ref{pf:sh-div:V:dens} (each $L_v$ is self-dual), whose product is
given by \ref{pf:sh-div:V:euler} and produces $L(n,\eta^n)^{-1}$ and the discriminant factors; the archimedean integral is given by
\ref{pf:sh-div:V:real}. On the left side, \ref{pf:sh-div:V:unf} expresses the mean through $\vol_D(W_1)$ and $\vol_D(W')$ (these are the
measures of fundamental domains by $\sfname{VolTie}(D)$) and the orbital integral of $\Phi_\infty$ of \ref{pf:sh-div:V:real}. Comparing
the two sides and using $[F:\Q]=f$ (the number of real places of $F$) gives \eqref{eq:pf:sh-div:vr}.
\end{proof}

\begin{remark}\label{pf:sh-div:rem:volratio}
The only hypothesis is the measurable one-vector mean value formula $\sfname{SW1b2}$; the finiteness of the volumes enters through
$\sfname{VolTie}(D)$, which is a hypothesis of $\sfname{MassRatioUnderVolume}(D)$. The six statements
\ref{pf:sh-div:V:perp}--\ref{pf:sh-div:V:unf} are proved in the development without hypotheses; their proofs are not reproduced here.
\end{remark}

%==========================================================================================================================
\subsection{Theta series: rationality, integrality and Galois conjugates}\label{sec:pf:sh-div:theta}

\paragraph{Setting.} We use the objects of \cref{sec:pf:liftinj,sec:pf:onto3}. A \emph{lift setting} consists of a global datum
$D$ (ranks, host $E\supseteq K$, spaces), a family $A$ (with base type $\mu_0$ and base embedding $\sigma_0$ of the coefficient field
$\mathit{LL}$), a core $T$ (with its moduli data, levels and the base embedding $\sigma_0$), a real datum $R$ over the coefficient field
of $T$, a face $F$, weight families $\kappa=(\kappa_i)_{i<m}$ and vectors $v_0=(v_0(i,\sigma))$. The \emph{standard hypotheses} on a
lift setting are: the laws of the core; the level characters of $R$ are standard with base type $\mu_0$; the group, family, model
and theta-pin conditions of $R$ (relative to $F$); and the tie of $(T,R,\kappa,v_0)$. For an index $i<m$ such that $T$ has the space
$W_c(i)$ at level $i+1$ and $\sigma_0$, the \emph{frame pin} asserts that the theta values of an orthonormal lattice basis are pinned
by a phase (\cref{sec:pf:liftinj}), and $\sfname{VStd}(D,i)$ that the space $V_c(i)$ is standard. A \emph{lowest weight vector}
$u=(u_\rho)$ is a family of vectors of the lowest $K$-types of the real models whose components are the $v_0(i,\sigma_0)_\rho$. For
such data the development defines (\cite{LeanSource2026}): the theta function $\theta_\omega(s,g,\phi)$ of the Weil representation
$\omega$ attached to the base characters, for $s$ in the finite adelic group of $V_c(i)$ and $g$ in the adelic group of $W_c(i)$;
the harmonic vectors $\mathrm{harm}(u,w,f)$ attached to a level datum $f$ and a rational vector $w$; \emph{covers} $C$ of the moduli
datum of $T$ at level $i+1$ (with a coefficient field $C.F$, an embedding $C.\sigma_2$, a level $C.l_2$ with auxiliary part
$C.p_d$, integral forms and a notion of rational function, $\mathrm{IsRat}$); and \emph{good} covers (uniformised, pinned at all
dominant weights, of auxiliary split squarefree level, small, and with field containing the ray class field of conductor
$M\,C.p_d$, $M$ the level of $T$).

\paragraph{The three series statements.}
\begin{itemize}
\item $\sfname{UntwRatS}$ (rationality of untwisted series): for every lift setting with the standard hypotheses, every index $i$
with frame pin and $\sfname{VStd}(D,i)$, every family tie $(\mathrm{gram}_R,c_{0R},c_{0L})$ and every lowest weight vector $u$, there
is $\mathrm{ph}\in\C$ with $|\mathrm{ph}|=1$ such that for every good cover $C$, every $s$, every level datum $f$ at level $M\,C.p_d$
with values in $C.\sigma_2(C.F)$ and every rational vector $w$ (for $c_{0L}(i)$ and $C.\sigma_2(C.F)$), the function
$g\mapsto\mathrm{ph}\cdot\theta_\omega(s,g,\mathrm{harm}(u,w,f))$ is rational on $C$.
\item $\sfname{SeriesRatS}$ (the statement (RS) of \cref{sec:pf:onto3}): the same, with the conclusion strengthened to: this
function is rational, and for every twist $(\delta,\delta_c)$ of conductor $C.p_d$ the products $dd\cdot\mathrm{ph}\,\theta_\omega$ with
every function $dd$ with the $\mu$-values of $\delta_c$, and the functions $\mathrm{ph}\,\theta_{\omega'}$ for every Weil datum $\omega'$
of the characters twisted by $\delta_c$, are rational on $C$.
\item $\sfname{SerInt}$ (integrality): for every lift setting with the standard hypotheses and every index $i$ with
$\sfname{VStd}(D,i)$ and frame pin, (a) for every lowest weight vector $u$ and every coefficient map $c_0$ (the base change of
$c_{0L}(i)$ for some family tie) there is a phase $\mathrm{ph}$ such that for every prime $l$, every $s$ and every finitely generated
$\Z$-lattice $\mathrm{Lat}$ of vectors there is $c\in\N$ with: for every tower cover $C$ at $l$ pinned at the dominant weights and
with large field, every set $S$ of primes of $C.F$ above $l$, every twist $\xi$ of conductor $C.p_d$, every level vector $\Phi$ with
$S$-integral values and associated $f$, and every $v\in\mathrm{Lat}$, there is an $S$-integral form $\phi$ on $C$ with
\[
\phi(g)=l^c\,\mathrm{ph}\,\xi(\det g)\,\theta(u,s,f,c_0(v))(g)\qquad\text{for all }g;
\]
and (b) ($\sfname{DescIntAt}$) every $S$-integral form on a tower cover at $l$ which is pinned at the dominant weight, unramified at
$S$ and invariant under the level group is an $S$-integral linear combination of integral forms of the base.
\end{itemize}
Finally $\sfname{GlobalArtinReciprocity}$ is the existence, for every number field $E$ and algebraic closure $\Omega$, of a
reciprocity map $\A_E^\times\to\Gal(\Omega/E)^{\mathrm{ab}}$ satisfying the axioms of global class field theory (open quotient map,
kernel the closure of $E^\times$ times the connected component, local compatibility with inertia), and the \emph{ray class field
statement} asserts that for every number field $K$ and non-zero ideal $\mathfrak q$ there is an abelian extension $L/K$, unramified
outside $\mathfrak q$ and generated by the Frobenius elements of the primes prime to $\mathfrak q$, such that an ideal $I$ prime to
$\mathfrak q$ has trivial Artin symbol if and only if it is principal in the ray modulo $\mathfrak q$.

\begin{theorem}[Rationality and integrality of theta series]\label{pf:sh-div:series}
Global Artin reciprocity implies $\sfname{SeriesRatS}$ and $\sfname{SerInt}$; in fact $\sfname{SerInt}$ and $\sfname{UntwRatS}$ hold
unconditionally.
\formallabel{GlobalTheory.seriesRatS_and_serInt_of_globalArtin}
\formallabel{GlobalTheory.T.thetaSeriesRationalOnGoodCovers_of_untwisted}
\formallabel{GlobalTheory.SerStdTwist.seriesRatS_of_untwRatS}
\formallabel{GlobalTheory.T.thetaSeries_integral_uniform_prime_power}
\formallabel{GlobalTheory.serInt_outright}
\end{theorem}

\begin{proof}
\emph{Step 1: the ray class field.} Global Artin reciprocity implies the ray class field statement
(\lref{GlobalTheory.Published.GlobalArtinReciprocity.rayClassFieldSt}).

\emph{Step 2: the key bound.} The common source of rationality and integrality is the following statement
(\lref{GlobalTheory.LevelSeries.series_bound}). Let $(V,W,\omega,\Lambda_V,\Lambda_W,f_0)$ be a \emph{vacuum pair} (hermitian
spaces with Weil representation, lattices and a vacuum vector $f_0$, with the standard properties of the theta data) with $W$ of
odd rank $r$, $\chi$ a character datum of level $M_0\ne0$, and $l$ a prime. For $s$ in the finite adelic group of $V$ there is
$c\in\N$ such that for every level $l_2$ with $l_2.M=M_0p_d$ ($p_d$ prime to $M_0$ and $l$, $l_2.M$ prime to the bad primes, $l_2$ good
at $l$), every coefficient field with a good embedding and every compatible dictionary, every prime $\lambda\mid l$ of the
coefficient field with valuation ring $\OO$, every $\OO$-valued level function $\Phi$ and theta datum $\phi$, and every harmonic
polynomial $Q$ with coefficients in $\OO$, there are $\beta_k\in\OO$ and holomorphic sections $f_k$ of the automorphic line bundle on
the moduli space, $\lambda$-integral at all points, with
\[
l^c\,\theta_\omega\bigl(s,g,\mathrm{std}(\phi,Q)\bigr)=\sum_k\beta_k\cdot\mathrm{Lie}\cdot c(g)\,\mu(g)\,f_k(g)\qquad\text{for all }g,
\]
where $\mathrm{Lie}$ is an explicit positive constant and $c,\mu$ are the dictionary constants. Its formal proof (about $10\,000$ lines)
bounds the level of the vacuum pair uniformly, moves the series by a central flip to half level and by translates with bounded
denominators, lets the real rational points act through a harmonic substitution, and identifies the theta series on covers with
short theta series of forms, whose values are bounded at all points of good reduction
(\lref{GlobalTheory.T.exists_forms_theta_short_series_identity}, \lref{GlobalTheory.T.exists_centralFlip_thetaTranslate_halfLevel}, \lref{GlobalTheory.T.exists_levelSeries_translate_boundedDenominator}, and \lref{GlobalTheory.T.real_rational_point_acts_harmonic_substitution});
it is not reproduced here.

\emph{Step 3: untwisted rationality.} $\sfname{UntwRatS}$ holds (\lref{GlobalTheory.T.untwisted_theta_rational_on_good_covers}).
Under the standard hypotheses, the theta data of $(V_c(i),W_c(i))$ form a vacuum pair (\lref{GlobalTheory.AdTh.RecPair.isPair_record_ofFive})
with $W_c(i)$ of odd rank $r$ ($k$ and $m$ are even), and the reference character datum has level $M=g_0g_1$, prime to the bad
primes. The core constants agree with the vacuum identity up to a phase $\mathrm{ph}_1$
(\lref{GlobalTheory.T.coreConstants_eq_vacuumIdentity_up_to_phase}), and the harmonic map sends the extremal vector to
$\kappa_a$ times the letters, $|\kappa_a|=1$ (\lref{GlobalTheory.T.harmonicMap_extremalVector_eq_phase_smul_letters}); the phase of
$\sfname{UntwRatS}$ is $\mathrm{ph}=(\kappa_a\mathrm{ph}_1)^{-1}$. Given a good cover $C$, choose a prime $l>C.p_d$; the cover is a
dictionary compatible with the frames and its constants are the vacuum identity constants
(\lref{GlobalTheory.T.uniformisedCover_isDictionary_frameCompatible} and \lref{GlobalTheory.T.vacuumIdentityConstants_cover_of_dominantPins}),
and Step~2 at $l$ (with $\OO$ the valuation ring at a prime of the coefficient field of $C$ above $l$) expresses
$l^c\,\mathrm{ph}\,\theta_\omega(s,\cdot,\mathrm{harm}(u,w,f))$ as a combination, with coefficients in the field of $C$, of rational
sections; this combination is rational on $C$, and so is its product with the element $l^{-c}$ of the field of $C$.

\emph{Step 4: twisting.} The ray class field statement and $\sfname{UntwRatS}$ imply $\sfname{SeriesRatS}$. Let $\mathrm{ph}$ be the
phase of $\sfname{UntwRatS}$. For a good cover $C$, the untwisted function is rational. For a twist $(\delta,\delta_c)$ and $dd$ with
the $\mu$-values of $\delta_c$, the product of a rational function on $C$ with such a determinant twist is rational: the values of
$\delta_c$ lie in the ray class field of conductor $M\,C.p_d$, which is contained in the field of the good cover
(\lref{GlobalTheory.T.rationalForms_detTwist_mul_rational}, which uses the ray class field statement). For a Weil datum
$\omega'$ of the twisted characters, $\theta_{\omega'}=c\cdot dd\cdot\theta_\omega$ for a constant $c$ and a function $dd$ with the
$\mu$-values of $\delta_c$ (\lref{GlobalTheory.WeilTwistSeries.theta_weil_twist}), and a constant multiple of a rational function
by a constant of the field of $C$ is rational.

\emph{Step 5: integrality.} $\sfname{SerInt}$ holds: (a) follows from Step~2 applied at the prime $l$, with $c$ uniform over the
lattice $\mathrm{Lat}$ (\lref{GlobalTheory.LevelSeries.serIntAt_record}), and (b) is the unramified descent of integral forms on
tower covers (\lref{GlobalTheory.T.unramified_descent_integral_forms_tower_cover}: integral forms on a tower cover unramified at
$S$ and invariant under the level group are combinations of base forms with $S$-integral coefficients, by rational descent of cover
forms to the base and descent at one place at a time).
\end{proof}

\paragraph{Galois conjugates of rational theta lifts.}
For a lift setting, an index $i$ and a phase $\mathrm{ph}$, a vector $y\in\mathit{LL}^{\,n_W(i)}$ is a \emph{rational lift} of $x$ (a vector
of the family at level $i$ and base type $\mu_0$) if $j_{\sigma_0}(y)=\mathrm{ph}\cdot\theta_G(j_V(x))$, where $\theta_G$ is the global theta
lift of $R$ at $(i,\mu_0)$ and $j_\sigma$ the embedding of the face at $\sigma$. An embedding $\sigma:\mathit{LL}\to\C$ is \emph{over $K$} if
it agrees with the base embedding on $K\subseteq E\subseteq\mathit{LL}$. The \emph{conjugation property} $\mathrm{Conj}(\sigma)$ asserts:
the hook sums at the label of $\sigma$ equal those at the identity; conjugation data exist; and there are a phase function $u$, a
unitary operator $U$ of the space of the face at $\sigma$, and a central shift $(z,\zeta)$ ($u=1$ if $\sigma=\sigma_0$) such that for every
conjugation datum the conjugate hypotheses hold and there is $\kappa$ with $|\kappa|=1$ and
$U(j_\sigma(y))=\kappa\,g_W(\sigma,z)\bigl(\mathrm{th}(j(x))\bigr)$ for all rational lifts $y$ of $x$, i.e.\ the $\sigma$-conjugate of the
rational lift is, up to a unimodular constant and a unitary operator, the theta lift of the conjugate datum shifted by $z$.
$\sfname{GalOut}(i,\mathrm{ph})$ is the conjunction of: (G1) every $x$ has exactly one rational lift; (G2) for every family tie and
every $\sigma$ over $K$, $\mathrm{Conj}(\sigma)$; (G3) if $\sigma$ is not over $K$, then $\bar\sigma$ is over $K$ and
$\langle j_\sigma y,j_\sigma y'\rangle=\overline{\langle j_{\bar\sigma}y,j_{\bar\sigma}y'\rangle}$ for all $y,y'$.

$\sfname{ConjRatLift5}$ is the statement: for every lift setting with one-map data $O=(O_i)$ at the base type, and every phase
$\mathrm{ph}$, if the five core conditions hold (standard level characters, group, family, model and theta-pin conditions), the base
character is standard of sign $-1$, there is a tie $\mathrm{fw}$ with its bar tie, the extension condition and the uniformity of $R$
hold and the one-map pins hold at every index, then $\sfname{GalOut}(i,\mathrm{ph})$ holds at every index $i$ at which $O_i$ is a
one-map datum for $(\kappa_i,v_0(i,\sigma_0))$ with phase $\mathrm{ph}$.

\begin{theorem}[Conjugates of rational theta lifts]\label{pf:sh-div:conj}
$\sfname{ConjRatLift5}$ holds.
\formallabel{GlobalTheory.T.conjugatesOfRationalThetaLift_from_characterisation}
\formallabel{GlobalTheory.ConjRatLiftFive.conjRatLift5}
\end{theorem}

\begin{proof}
Fix the data and an index $i$ as in the statement. From the core conditions, the tie and the pins one obtains the evaluation
readings $j_{\mathrm{ev}}$ and $j_{\overline{\mathrm{ev}}}$ of the tie (\lref{GlobalTheory.ConjReading.jEvExists_and_jEvBar}); in
particular $\mathrm{fw}$ evaluates, and the evaluation at unit vectors is injective
(\lref{GlobalTheory.EvalAtUnitVectorFinal.evInj_of_isModels_of_tie}). The embedding $j_{\sigma_0}$ of the face at the base
embedding is injective, since the family of forms at $\sigma_0$ is linearly independent
(\lref{GlobalTheory.GalOne.inj_of_indep}).

(G1) Choose representatives for the pins at $i$. The base lift construction gives, for every $x$, a vector $y$ with coordinates in
$\mathit{LL}$ and $j_{\sigma_0}(y)=\mathrm{ph}\cdot\theta_G(j_V(x))$ (\lref{GlobalTheory.ConjRatLiftFive.base_lift5}); it is unique because
$j_{\sigma_0}$ is injective.

(G2) Let $\sigma$ be over $K$. If $\sigma=\sigma_0$, $\mathrm{Conj}(\sigma_0)$ is deduced from the core conditions, the tie, the
injectivity of the evaluation, the family tie and the one-map datum at $i$ (\lref{GlobalTheory.ConjRatLiftFive.conjAt_base5}). If
$\sigma\ne\sigma_0$, $\mathrm{Conj}(\sigma)$ is the main statement \lref{GlobalTheory.ConjRatLiftFive.conjugateRationalLift_at_nonBaseEmbedding},
which carries most of the formal proof. Its ingredients are: the conjugation of one-map data, which identifies the $\sigma$-conjugate
of the rational lift with the theta lift of the conjugate datum (\lref{GlobalTheory.ConjRatLiftFive.conj_lift5}, resting on
\lref{GlobalTheory.ConjRatLiftFive.oneMapConj5}); the conjugate lift obtained from the characterisation clauses
(\lref{GlobalTheory.T.conjugateLift_of_characterisationClauses}); and the conjugate extremal vectors provided by moved weights
(\lref{GlobalTheory.MovedWeights.exists_io_conj_ofFive}). Three further statements of the development are used in the formal proof
of the theorem: the uniqueness of conjugate lifts together with the base datum
(\lref{GlobalTheory.T.conjugateLift_unique_and_baseDatum}), the conjugate characters and embeddings
(\lref{GlobalTheory.T.conjugateCharacters_embeddings_of_characterisation}), and the central shifts and phases of the conjugate readings
(\lref{GlobalTheory.T.centralShift_phase_conjReadings_fiveClauses}).

(G3) If $\sigma$ is not over $K$, then $\bar\sigma$ is over $K$, since $\sigma$ and $\bar\sigma$ induce the two embeddings of the imaginary
quadratic field $K$ (\lref{GlobalTheory.GalLift.labels_overK}); and the identity of inner products is the third clause of
\lref{GlobalTheory.T.centralShift_phase_conjReadings_fiveClauses}, applied with the readings $j_{\mathrm{ev}}$,
$j_{\overline{\mathrm{ev}}}$ and the bar tie, using that $\bar\sigma$ is good exactly when it is over $K$
(\lref{GlobalTheory.GalLift.overK_iff_good_ofFive}).
\end{proof}

\begin{remark}\label{pf:sh-div:rem:conj}
The case $\sigma\ne\sigma_0$ of (G2) carries almost all of the formal proof (about $9\,800$ of its $11\,900$ lines); its ingredients are
listed above and their proofs are not reproduced here.
\end{remark}

\paragraph{The class sum span.} For a lift setting and a one-map datum $O$ at index $i$ (with coefficient field $c'$), the
\emph{cusp forms over $O$} of weight $\kappa_i$ form a $c'.F$-vector space; a \emph{class sum} is a cusp form over $O$ which is the
theta class sum attached to a point $x$ of $O$ (\cite{LeanSource2026}, \cref{sec:pf:onto3}). $\sfname{ClassSumsSpan}(O,\kappa_i)$ asserts
that the cusp forms over $O$ of weight $\kappa_i$ lie in the $c'.F$-span of the class sums.

\begin{theorem}[Class sums span the cusp forms]\label{pf:sh-div:classsum}
Let a lift setting with one-map data $O=(O_i)$ at the base type be given, with $8fr+4\le n_V$ (the rank of the spaces $V$), the core
conditions $\sfname{Core}(R,F)$, a tie $\mathrm{fw}$ and the extension condition. Then for every index $i$ at which the one-map pins
hold (and the one-map row condition holds), $\sfname{ClassSumsSpan}(O_i,\kappa_i)$ holds.
\formallabel{GlobalTheory.ClassSumSpanCore.classSumsSpan_core}
\formallabel{GlobalTheory.ClassSumSpanCore.classSumsSpan_of_pins_core}
\end{theorem}

\begin{proof}
Let $f$ be a cusp form over $O=O_i$ of weight $\kappa_i$. Let $\mathrm{rd}$ be the map sending a form to its unfolded section on the
moduli datum of the core, and $\mathrm{at}$ the evaluation of such sections along the dictionary of $O$; the composite
$\mathrm{at}\circ\mathrm{rd}$ is injective on the cusp forms over $O$ (\lref{GlobalTheory.ClassSumSpan.atDict_indep}).
Hence it suffices to show that $\mathrm{at}(\mathrm{rd}(f))$ lies in the span of the $\mathrm{at}(\mathrm{rd}(g))$, $g$ a class sum.
\begin{enumerate}[label=(\arabic*)]
\item \emph{Raw series.} $\mathrm{rd}(f)$ lies in the span of the \emph{raw series} of the core at the points of a fixed good point $x_1$
of $O$, moved by unitary transformations of the datum (\lref{GlobalTheory.ClassSumSpan.rdU_mem_span_raw}; the formal proof moves
the datum by suitable unitary matrices, \lref{GlobalTheory.DatumMove.rawSeries_move_datum}); here the rank condition $8fr+4\le n_V$
is used.
\item \emph{Splitting the move.} Each raw series is, after a split move, a non-zero constant $z$ times a function $t$ whose value at
every dictionary point is a raw value at a point (\lref{GlobalTheory.ClassSumSpan.exists_split_move}).
\item \emph{Kernel values.} At each dictionary point, the raw value equals a theta kernel value at a finite point, multiplied by
explicit non-zero constants: the frame constants of the characterisation core, the Lie factor of the weight, the $\mu$-values of the
dictionary and the phase of the Fock coefficient map of the extremal vector
(\lref{GlobalTheory.ClassSumSpanCore.kernel_eq_shortVal}; it uses the identification of the one-class theta kernel with the
classical short theta series, \lref{GlobalTheory.T.thetaKernel_oneClass_eq_classicalShortSeries}, and
\lref{GlobalTheory.T.fockCoefficientMap_extremal_eq_phase_smul_letters}).
\item \emph{Rational form lifts.} The theta kernel at a finite point lies in the span of the lifts of rational forms
(\lref{GlobalTheory.T.thetaKernel_finitePoint_mem_span_rationalFormLifts}), and every such lift is, up to the constant
$u$ of modulus $1$, a class sum: by the adelic lift identity the lift is the geometric class sum at every tuple
(\lref{GlobalTheory.T.exists_adelicLift_eq_geometric_classSum_all_tuples}), and the class sum rule provides a cusp form over $O$
which is a class sum at the corresponding point (\lref{GlobalTheory.ClassSumRule.exists_cusp_isThClsD_core}).
\end{enumerate}
Since all constants are non-zero, (1)--(4) give $\mathrm{at}(\mathrm{rd}(f))$ in the span of the $\mathrm{at}(\mathrm{rd}(g))$.
\end{proof}

\begin{remark}\label{pf:sh-div:rem:classsum}
The hypothesis that the one-map row condition holds at $i$ is not used, and of the one-map pins only the second component is used.
\end{remark}

\paragraph{Rational rows.} We use the step statements of \cref{sec:pf:rowprops,sec:pf:stageU}. For properties $p,q$ of column
witnesses $w$ over a base, $\sfname{StepW}(p,q)$ is the eventual statement (for all large parameters, at every context over the base
with forced data) that $p(w)$ implies $q(w)$ for every column witness $w$; it is transitive, and it holds whenever $p(w)\Rightarrow q(w)$
for all contexts and witnesses (\lref{GlobalTheory.Mem.Top.Bl.StepW.of_all}). $\sfname{Pinned}(w)$ is the existence of column pins,
and $\sfname{Start}(w)$ is $\sfname{Pinned}(w)$ together with the support and coherence conditions. For a property $P$ of data at a
witness, $\sfname{AtDat}(P)=\sfname{StepW}(\sfname{Pinned},Q_P)$, where $Q_P(w)$ asserts $P$ at every later stage $Z$, every field
condition, every set of places $\mathrm{Wt}$ and every datum $z$ for which the group is true at $Z$, the group condition at the stage,
the existence condition of the real datum and the geometric condition hold. $\sfname{RowRatW}=\sfname{StepW}(\sfname{Start},Q')$, where
$Q'(w)$ asserts the row test $\sfname{RowTest}$ of the complex datum of $z$ under the same hypotheses and, in addition, the
$\kappa$-definition, the core property, the group condition $\sfname{GIs2}$, $\mathrm{Wt}=W_1$ and the condition that the residue
characteristics of the places of $\mathrm{Wt}$ lie in the places of $T$.

\begin{proposition}[Rational rows from the row test]\label{pf:sh-div:rowrat}
$\sfname{AtDat}(\sfname{RowTest})$ implies $\sfname{RowRatW}$.
\formallabel{GlobalTheory.Mem.LiteralAutomorphy.rowRatW_of_rowTest}
\end{proposition}

\begin{proof}
Write $Q_{\mathrm{RT}}$ for $Q_P$ with $P=\sfname{RowTest}$. Since $\sfname{Start}$ contains $\sfname{Pinned}$, we have
$\sfname{StepW}(\sfname{Start},\sfname{Pinned})$. Composing with $\sfname{AtDat}(\sfname{RowTest})=\sfname{StepW}(\sfname{Pinned},Q_{\mathrm{RT}})$
gives $\sfname{StepW}(\sfname{Start},Q_{\mathrm{RT}})$. Finally $Q_{\mathrm{RT}}(w)$ implies $Q'(w)$, because $Q'$ asks for the same
conclusion under more hypotheses; by transitivity, $\sfname{StepW}(\sfname{Start},Q')$, which is $\sfname{RowRatW}$.
\end{proof}

\paragraph{Symmetric column forms.} For a base $B$ (\cref{sec:pf:open0}), a pre-datum $X$ over it and a column input $I$ with the
condition $I.\mathrm{Ok}$ and the pins, laws $l$ of the columns of $I$ and later stages $Z$, the record $I.\mathrm{rec}_7(l,Z)$ is a
global datum $D$; \emph{symmetric tied column forms} for $D$ are the hermitian column data with, for every word $u$, tied symmetric forms
$C.\mathrm{forms}(u)$, and $\chi\mathrm{Fin}_0$ for such forms asserts that the level character of the step-zero class has finite order
and takes values in $\sigma_0(\mathit{LL})$.

\begin{proposition}[Symmetric column forms with finite step-zero characters]\label{pf:sh-div:symcol}
In this situation there are symmetric tied column forms $C$ for $I.\mathrm{rec}_7(l,Z)$ such that $\chi\mathrm{Fin}_0(C.\mathrm{forms}(u))$
holds for every word $u$.
\formallabel{GlobalTheory.T.exists_symmetricColumnForms_stepZeroFiniteCharacter}
\formallabel{GlobalTheory.Mem.Col2.TC.colsTiedSym_chiFin0}
\formallabel{GlobalTheory.formsTiedSym_chi_zero_finite}
\end{proposition}

\begin{proof}
The standard column comparisons of the true columns give symmetric tied column forms $C$
(\lref{GlobalTheory.T.exists_standardColumnComparisons_trueColumns}). For each word $u$, the step-zero level character of
$C.\mathrm{forms}(u)$ is the character of a unipotent-type step level at the level primes $g_0,g_1$; such a character has finite
order, and its values lie in $s(K)$ for every field $K$ with an embedding $s:K\to\C$ which contains primitive $g_i$-th roots of unity
for $i=0,1$ (\lref{GlobalTheory.T.unipotentTypeLevel_character_finiteOrder_rootsOfUnity}). The coefficient field $\mathit{LL}$ of the
record contains primitive $g_i$-th (and $(g_i-1)$-th) roots of unity for $i=0,1$
(\lref{GlobalTheory.T.exists_primitiveRoots_coeffField_trueColumns}); apply this with $K=\mathit{LL}$ and $s=\sigma_0$. (The
$(g_i-1)$-th roots of unity are not used.)
\end{proof}

%==========================================================================================================================
\subsection{Host frames, twisted places, numerical rows and the cut inequality}\label{sec:pf:sh-div:frames}

\paragraph{Contexts and eventual statements.} A \emph{context} $K$ (\cref{sec:pf:baselat1,sec:pf:bundles}) consists of a scale $s$
(with exponent $\eps$, scale parameter $x$ and $L=x^\eps$), standard constants $C$, real parameters $c_D,c_h$ and a pre-base over which
it lies; a \emph{stage tuple} $X$ over a pre-base consists of a base and its column, row, divisibility, $S$-, $N$- and $GJ$-stages, and
determines a global datum $D=X.D$. For a property $P$ of contexts, $\shdivEv(P)$ is the statement: there is $\eps_0>0$ such that for
all $\eps\le\eps_0$ with $\eps<\frac1{52}$ and all $C_8\ge1$ there is a transported constant $K_0$ such that for all standard constants
$C$ with $C_{pde}=C_8$ which are large with respect to $K_0$ there is $c_1$ such that for all $c_D\ge c_1$ there is $c_2$ such that for all
$c_h\ge c_2$ there is $x_0$ such that $P(K)$ holds for every context $K$ with constants $C$, parameters $c_D,c_h$, exponent $\eps$ and
$x\ge x_0$. If for every $\eps\in(0,\frac1{52})$ there are $y,x_0$ such that $P(K)$ holds whenever $K$ has exponent $\eps$, $c_D,c_h\ge y$ and
$x\ge x_0$, then $\shdivEv(P)$: take $\eps_0=1$, any $K_0$, $c_1=c_2=y$ and the given $x_0$.

For a global datum $D$ with ranks $m,k,t$, $r$ and $N(m)$, and Harish-Chandra parameters $(b_j,e_j)_{j<m}$, $(v_i)_{i<k-1}$ at each
$\tau$, with associated sequence $\lambda_n$ ($\lambda_n=v_n$ for $n<k-1$ and $\lambda_n=-b_{n-k+1}$ for $k-1\le n<r$), the \emph{host
hypothesis} $\sfname{HostHyp}(D)$ asserts: the parameters interlace; $e_j-b_j\ge\frac{t+3}2$ and $b_{j+1}-e_j\ge\frac{t+3}2$; $\lambda_n-\lambda_{n+1}\ge
100r^4$ for $n+1<r$; the cloud gap $\lambda_{k-2}-\lambda_{k-1}\ge100r^4$; and the node gaps of $D$ are at least $4N(m)^3+2t+6$ and at
least $256r^4$. A hermitian space is \emph{$F$-diagonal} if its Gram matrix in its frame is diagonal with entries in $F$, and it has a
\emph{standard lattice} if its lattice is the set of vectors with integral coordinates.

\begin{theorem}[Host hypotheses and frames, eventually]\label{pf:sh-div:hostframes}
$\shdivEv$ holds for the following property of contexts $K$: for every stage tuple $X$ over the pre-base of $K$, with $D=X.D$,
\begin{enumerate}[label=(\roman*)]
\item $\sfname{HostHyp}(D)$, and $F$ is the maximal real subfield of $E$;
\item for all $i$ the spaces $W_c(i)$ and $V_c(i)$ are $F$-diagonal and $V_c(i)$ is positive definite;
\item the spaces $Y_j$, $V_j$, $X'$ and $W_j$ ($1\le j\le m$) have standard lattices;
\item the frames of $V_c(i)$ and of $X'$ are orthonormal and span their lattices over $\OO_E$.
\end{enumerate}
\formallabel{GlobalTheory.Mem.HostFrames.hostHyp_frames_ev}
\formallabel{GlobalTheory.Mem.HostFrames.hostHyp_of_gaps}\formallabel{GlobalTheory.Mem.HostFrames.ev_of_large}
\end{theorem}

\begin{proof}
By the criterion above it suffices, for $\eps\in(0,\frac1{52})$, to find $y,x_0$. Let $y=402\cdot4+2$ and let $x_0$ be the maximum of
the thresholds of the following two size statements (\lref{GlobalTheory.Mem.HostFrames.size_le} and \lref{GlobalTheory.Mem.HostFrames.thin_size}):
for scales with exponent $\eps$ and $x\ge x_0$ and ranks with the rank valuations of the scale, $256\,N(m)^4\le\lceil x^{y'}\rceil$ for all
$y'\ge y$, and $100r^4+1\le w''(n)$. Let $K$ be a context with $c_D,c_h\ge y$ and $x\ge x_0$, and $X$ a stage tuple over its pre-base. Since
$X$ lies over the pre-base of $K$, its ranks satisfy the rank valuations of the scale and its Harish-Chandra parameters are prescribed
by the context (\lref{GlobalTheory.Mem.HostFrames.val_of_over}): for every word $u$ and every $\tau$, $e_j-b_j=b_{j+1}-e_j=h:=\lceil
x^{c_h}\rceil$ and $v_i=(k-2-i)\cdot4\Delta$ with $\Delta=\lceil x^{c_D}\rceil$; moreover, since $100r^4+1\le w''(n)$, the first node satisfies $b_0\ge100r^4$
(\lref{GlobalTheory.Mem.HostFrames.first_node_over}). By the size statements, $h\ge256N(m)^4$ and $\Delta\ge256N(m)^4$.

\emph{(i)} We show that these conditions imply $\sfname{HostHyp}(D)$. Note $r\le N(m)$ and $t\le N(m)$. The interlacing and the two gap
conditions $\ge\frac{t+3}2$ are part of the data of Harish-Chandra parameters. Consecutive $b_j$ differ by $2h$. For the gaps of the
$\lambda_n$: if $n+1<k-1$, $\lambda_n-\lambda_{n+1}=4\Delta\ge1024N(m)^4\ge100r^4$; if $n+1=k-1$, $\lambda_{k-2}-\lambda_{k-1}=0+b_0\ge100r^4$ (this
is also the cloud gap); if $n\ge k-1$, $\lambda_n-\lambda_{n+1}=b_{j+1}-b_j=2h\ge100r^4$. The node gaps are differences of consecutive $b_j$,
i.e.\ $2h\ge512N(m)^4$, which exceeds both $4N(m)^3+2t+6$ and $256r^4$. The statement that $F$ is the maximal real subfield of $E$ is
part of the data of a stage over a pre-base (\lref{GlobalTheory.Mem.HostFrames.frames_of_over}).

\emph{(ii)--(iv)} Over the pre-base, the spaces $W_c(i),V_c(i),Y_j,V_j,W_j,X'$ have \emph{sign frames}: frames $(e_a)$ with
$\langle e_a,e_b\rangle=\pm\delta_{ab}$ spanning the lattice; $V_c(i)$ is positive definite and the frames of $V_c(i)$ and $X'$ are
orthonormal (\lref{GlobalTheory.Mem.HostFrames.frames_of_over}). A space with a sign frame is $F$-diagonal (its Gram matrix is
diagonal with entries $\pm1$). The spaces of the pre-base are diagonal in the sense of the development, which gives the standard
lattices (\lref{GlobalTheory.Mem.HostFrames.diagSpaces_of_over} and \lref{GlobalTheory.Mem.HostFrames.stdLat_diag}).
\end{proof}

\paragraph{Twisted places.} Let $D$ be a global datum with coefficient field $\mathit{LL}$, its subfield $\mathit{LL}_0$, its places
$w$ with types $\mu(w)$ and base type $\mu_0$, and twisting data $\mathrm{ntw}$: a set of twisted indices $i$ with local rings $\OO_{0,i}$
and embeddings $\mathrm{emb}_{0,i}:\mathit{LL}_0\to\mathrm{Frac}(\OO_{0,i})$. An index $i$ is a \emph{non-zero index place} if there is a place
$w'$ of $\mathit{LL}$ with $\mu(w')\ne\mu_0$ such that for $x\in\OO_{\mathit{LL}_0}$, $\mathrm{emb}_{0,i}(x)$ lies in the maximal ideal of
$\OO_{0,i}$ exactly when $v_{w'}(x)>0$ (i.e.\ $i$ \emph{lies under} $w'$). A \emph{choice} $c$ for $D$ assigns to every place $w$ a
twisted index $c.\mathrm{tw}(w)$ and a side $c.\mathrm{sd}(w)\in\{\mathrm{true},\mathrm{false}\}$; it is \emph{tied} to $D$
($\sfname{NChTie}$) if it is compatible with the local data of $D$ \cite{LeanSource2026}.

\begin{proposition}[Twisted place indices]\label{pf:sh-div:twisted}
Let $c$ be a choice tied to $D$, with $\mathit{LL}_0$ Galois over $\Q$. For every place $w$ with $\mu(w)\ne\mu_0$, the index
$c.\mathrm{tw}(w)$ is a non-zero index place.
\formallabel{GlobalTheory.Mem.Top.twistedPlaceIndex_mem_firstPair_of_compatibleChoice}
\end{proposition}

\begin{proof}
Let $w'=w$ if $c.\mathrm{sd}(w)$ is true and $w'=\bar w$ (the conjugate place) otherwise. Then $\mu(w')=\mu(w)\ne\mu_0$, since the type is
invariant under conjugation of places. Since $\mathit{LL}_0$ is Galois and $c$ is tied to $D$, the twisted index $c.\mathrm{tw}(w)$ lies
under $w'$ (\lref{GlobalTheory.TwistGoodCompare.liesUnder_of_tie}: an element $x$ of $\OO_{\mathit{LL}_0}$ is integral over
$\OO_{0,i}$, so $\mathrm{emb}_{0,i}(x)$ lies in $\OO_{0,i}$, and the tie identifies its non-units with the $x$ of positive valuation at
$w'$).
\end{proof}

\paragraph{Numerical rows.} For a property $P(C,D,L)$ of standard constants $C$, global data $D$ and size data $L$, $\sfname{SzRow}(P)$
asserts: for all $\eps\in(0,\frac1{52})$ and standard $C$ there is $c_1$ such that for all $c_D\ge c_1$ there is $c_2$ such that for all
$c_h\ge c_2$ there is $x_0$ such that $P(C,D,L)$ holds for all $D$ and size data $L$ with exponent $\eps$, scale $x\ge x_0$ and the
valuation conditions $\sfname{Val}(C,c_D,c_h,L)$ (prescribed ranks, constants and Harish-Chandra parameters). $\sfname{SzRow}$ is closed
under conjunction, and $\sfname{SzRow}(P)$ implies $\sfname{SzRow}(Q)$ if $P\Rightarrow Q$ under $\sfname{Val}$ for standard $C$
(\lref{GlobalTheory.EX.szRow_and_imp}). The \emph{host arithmetic} $\sfname{HostArith}(H)$ of a host asserts: $g_i=\pi_i\bar\pi_i$ with
$\pi_i\in\OO_E\cap K$, $g_i\equiv1\bmod Q_Sqd$, $p_{de}$ and $p'_{de}$ split completely in $E$, and all level primes have class
$\mathrm{splitK}$. The \emph{numerical record} $\sfname{NumAt}(C,L)$ consists of a parameter choice $P$ for $L$ with the constants of $C$
and $C^n_W(P)=C^n_W$, a second parameter record $T$ with the constants $(C_P,e_P,C_p,C_b,c_{TO},C_{hT},C_{bN})$ of $C$ and $e_P=C.e_P$, the
size conditions $\sfname{SizesAt}(L)$, the bound $\rho\le\cdots$ on the band, the identities for $C_\Gamma$, $C_c$, $C_t$, the second
constant, $\Lambda_G\ge x^\eps$, the bound for $b_T$, the top condition $c_{cK}$, $g_j\le e^{9L}$, $p_{de},p'_{de}<X'$, the log-five
depth, the identity for $L_w$ (in terms of $C_t$), the independence of the cloud parameters from $\tau$, and the windows for
$n,m,k$. Finally $\sfname{ChTT}(L)$ asserts $\frac12\delta_1n\le|T_T|\le2\delta_1n$ for the set $T_T$ of the host, and
$\sfname{ChCode}(L,3,C_F)$ asserts $3\Gamma_L^2x^{C_F}\le\#\{\text{words}\}$.

\begin{proposition}[Numerical data of the members]\label{pf:sh-div:szrow}
$\sfname{SzRow}$ holds for the property: if $\sfname{HostArith}(D.\mathrm{host})$, then a numerical record $\sfname{NumAt}(C,L)$ exists,
$\sfname{ChTT}(L)$ and $\sfname{ChCode}(L,3,C_F)$.
\formallabel{GlobalTheory.EX.szRow_numAt}
\end{proposition}

\begin{proof}
By closure under conjunction, the following four rows hold jointly:
\begin{enumerate}[label=(\alph*)]
\item if $\sfname{HostArith}$, a parameter choice $P$ with the constants of $C$ and $C^n_W(P)=C^n_W$ exists
(\lref{GlobalTheory.EX.explicitConstants_parameterConditions_large}: the explicit choice of the parameters satisfies all
parameter conditions at large scale);
\item for every such $P$ there is a second parameter record $T$ with the prescribed constants and $e_P$
(\lref{GlobalTheory.EX.szRow_blockLengthParams});
\item the size conditions $\sfname{SizesAt}(L)$ (\lref{GlobalTheory.EX.szRow_sizesAt});
\item the band bound $\rho_B\le x^{-2}$, $x^\eps\le\Lambda_G$, $g_j\le e^{9L}$, $p_{de},p'_{de}<X'$, $\sfname{ChTT}(L)$ and
$\sfname{ChCode}(L,3,C_F)$ (\lref{GlobalTheory.EX.szRow_band_lambda_primes}).
\end{enumerate}
By closure under implication it remains to assemble, under $\sfname{Val}(C,c_D,c_h,L)$ and $\sfname{HostArith}$, the record
$\sfname{NumAt}(C,L)$: its fields are $P$ and $T$ from (a), (b), the sizes from (c), the bounds from (d), and the identities for $C_\Gamma$,
$C_c$, $C_t$ (with $C_t=401$), the second constant, $b_T$, the log-five depth, the top condition and the windows, which are part of
$\sfname{Val}$. The identity for $L_w$ follows from the value $L_w=10\lceil405\log x\rceil+2\lceil30/\eps\rceil+22$ given by the host part of
$\sfname{Val}$ and $C_t=401$; the cloud parameters at $\tau$ and $\tau'$ coincide by the Harish-Chandra part of $\sfname{Val}$.
\end{proof}

\paragraph{The cut inequality for Weyl modules.} We use the notation of \cref{pf:onto2:weylexp}: for a dominant weight
$\nu=(\nu_0\ge\dots\ge\nu_{N-1})$, $\nabla_\Z(\nu)\supseteq\Delta_\Z(\nu)$ are the integral dual Weyl and Weyl modules inside the polynomial
ring $\Z[x_{ab}]$, $\nabla_\Q(\nu)=\nabla_\Z(\nu)\otimes\Q$, and $\shdivex_\ell(\nu)$ is the least $e$ such that $m\ell^e\nabla_\Z(\nu)\subseteq\Delta_\Z(\nu)$ for
some $m$ prime to $\ell$. The cut inequality $\sfname{ExCut}$ is the statement \eqref{eq:pf:onto2:excut} of \cref{sec:pf:onto2}: for all $N$,
all $\lambda:\N\to\N$ with $\lambda|_N$ antitone, every prime $\ell$, all $\mathrm{lo}\in\Z$, $w\in\N$ with $\mathrm{lo}\le\lambda(a')+N-1-a'\le
\mathrm{lo}+w$ for $a'<N$, and every $a\le N$,
\[
\shdivex_\ell(\lambda|_N)\le\shdivex_\ell(\lambda|_a)+\shdivex_\ell\bigl(\lambda(a+\cdot)|_{N-a}\bigr)+a(N-a)\lfloor\log_\ell w\rfloor .
\]
The space $\nabla_\Q(\nu)$ carries operators $e_{cd}$ ($c,d<N$; the action of the Lie algebra $\mathfrak{gl}_N$ by row derivations), a
highest weight vector $v_\nu$, and the linear form $\phi$ given by the Fock pairing with $v_\nu$. For $a+b=N$, the \emph{grid projector}
$\mathrm{grid}_K(T)$ of an endomorphism $T$ is an explicit ordered product, over the rows $i<a$, of row operators built from the
$\mathrm{ad}(e_{cd})$ and $K$, applied to $T$ \cite{LeanSource2026}.

The proof uses nine statements of the development, each proved there without hypotheses; we state them in the form used.
\begin{enumerate}[label=(W\arabic*)]
\item\label{pf:sh-div:W:levi} \emph{Levi projection} (\lref{GlobalTheory.T.dualWeylLeviProjectionDenominators_of_hypotheses}): for
$\nu$ antitone of length $p+q$ and $c_1,c_2$ with $c_1\nabla_\Z(\nu')\subseteq\Delta_\Z(\nu')$ and $c_2\nabla_\Z(\nu'')\subseteq\Delta_\Z(\nu'')$ for
the two blocks $\nu'=\nu|_p$, $\nu''=\nu(p+\cdot)|_q$, there is an endomorphism $T$ of $\nabla_\Q(\nu)$ which commutes with the $e_{cd}$ with
$c,d$ in the same block, maps $\nabla_\Z(\nu)$ into $(c_1c_2)^{-1}\Delta_\Z(\nu)$, kills the weight vectors whose weight has first-block sum
smaller than that of $\nu$, and satisfies $\phi(Tv_\nu)=\phi(v_\nu)$. (Its proof restricts $\nabla_\Z(\nu)$ to the Levi subgroup
$\GL_p\times\GL_q$, using the tensor structure of the restriction lattice, the block diagonal form of the restriction of the highest
weight vector and the row operators.)
\item\label{pf:sh-div:W:gl} \emph{$\mathfrak{gl}$-structure} (\lref{GlobalTheory.T.dualWeylModuleStructureGl_proof}): the $e_{cd}$ satisfy
the relations of $\mathfrak{gl}_N$, $\nabla_\Q(\nu)$ is spanned by weight vectors, $\nu$ is dominant for this action, and
$\nabla_\Z(\nu)$, $\Delta_\Z(\nu)$ are stable under the divided powers of the $e_{cd}$.
\item\label{pf:sh-div:W:bdd} \emph{Weight spread} (\lref{GlobalTheory.T.dualWeylWeightSpreadBound_proof}): if $\nu_i\le\nu_j+K$ for all
$i,j$, then $\nabla_\Q(\nu)$ is \emph{$K$-bounded}: every eigenvalue $n\in\Z$ of $e_{aa}-e_{bb}$ on $\nabla_\Q(\nu)$ satisfies $n\le K$.
\item\label{pf:sh-div:W:adj} \emph{Adjoint action} (\lref{GlobalTheory.T.adjoint_action_endomorphism_weight_bound}): then the operators
$\mathrm{ad}(e_{cd})$ on endomorphisms also satisfy the $\mathfrak{gl}$-relations and are $2K$-bounded.
\item\label{pf:sh-div:W:inv} \emph{Grid invariance} (\lref{GlobalTheory.T.grid_invariance_of_rank_one_two}, from invariance in ranks one and
two): if operators $D_{cd}$ satisfy the $\mathfrak{gl}$-relations and are $2K$-bounded, and $D_{cd}T=0$ whenever $c,d$ lie in the same
block, then $D_{cd}\,\mathrm{grid}_K(T)=0$ for all $c,d$.
\item\label{pf:sh-div:W:sc} \emph{Grid scalar} (\lref{GlobalTheory.T.projectorGridScalar_of_projectorScalar}, from the value of the
projector on a highest weight vector): for $T$ killing the weight vectors of smaller first-block sum and a highest weight vector $v$ of
weight $\lambda$ with $\lambda_i\le\lambda_j+K$,
$\phi(\mathrm{grid}_K(T)v)\cdot z=\phi(Tv)$ with $z=\prod_{i<p\le j<p+q}(\lambda_i-\lambda_j+j-i)$.
\item\label{pf:sh-div:W:schur} \emph{Schur's lemma} (\lref{GlobalTheory.T.dualWeyl_schur_endomorphism}): an endomorphism $S$ of
$\nabla_\Q(\nu)$ commuting with all $e_{cd}$ and with $\phi(Sv_\nu)=\phi(v_\nu)$ is the identity.
\item\label{pf:sh-div:W:int} \emph{Grid integrality} (\lref{GlobalTheory.T.projector_grid_denominator_bound}): if two lattices $L,M$ are
stable under the divided powers, $K+p+q\le w+1$, $\Lambda\in\N$ satisfies $(j+k+1)!\mid\Lambda\,j!\,k!$ for $j+k+1\le w$, and $cT(L)\subseteq M$,
then $c\Lambda^{pq}\,\mathrm{grid}_K(T)(L)\subseteq M$.
\item\label{pf:sh-div:W:beta} \emph{Beta denominators} (\lref{GlobalTheory.T.beta_integral_denominator_bound}): for every prime $\ell$
and $w$ there is $m$ prime to $\ell$ such that $(j+k+1)!\mid m\,\ell^{\lfloor\log_\ell w\rfloor}j!\,k!$ whenever $j+k+1\le w$.
\end{enumerate}

Two of these statements have short proofs, which we give.

\begin{lemma}[Schur's lemma for $\nabla_\Q(\nu)$, \ref{pf:sh-div:W:schur}]\label{pf:sh-div:schur}
Let $\nu$ be antitone of length $N$ and $S$ an endomorphism of $\nabla_\Q(\nu)$ commuting with all $e_{ab}$, $a,b<N$, with
$\phi(Sv_\nu)=\phi(v_\nu)$. Then $S=1$.
\formallabel{GlobalTheory.T.dualWeyl_schur_endomorphism}
\formallabel{GlobalTheory.LeviCut.nablaSchur}
\end{lemma}

\begin{proof}
For $a<b$ the raising operator $e_{ab}$ kills $v_\nu$ (\lref{Formalism.CohomologicalDescent.LR.der_hw_raise}), so
$e_{ab}(Sv_\nu)=S(e_{ab}v_\nu)=0$, and hence all divided powers of the row operators $(a,b)$, $a<b$, kill $Sv_\nu$. The vectors of
$\nabla_\Q(\nu)$ with this property form the line $\Q v_\nu$ (\lref{Formalism.GradedWeylDivisibility.hwLine}), so $Sv_\nu=cv_\nu$.
Since $\phi(v_\nu)$ is the Fock norm of $v_\nu$, which is positive (\lref{Formalism.CohomologicalDescent.LR.hw_norm_pos}), the
hypothesis gives $c=1$. Let $M=\ker(S-1)\subseteq\nabla_\Q(\nu)$. Since $S$ commutes with the $e_{ab}$, $M$ is stable under all $e_{ab}$,
hence under all row divided powers, and it contains $v_\nu$. A subspace of $\nabla_\Q(\nu)$ with these two properties is
$\nabla_\Q(\nu)$ itself (\lref{Formalism.GradedWeylDivisibility.closedEqNabla}). Hence $S=1$.
\end{proof}

\begin{lemma}[Beta denominators, \ref{pf:sh-div:W:beta}]\label{pf:sh-div:betaden}
Let $\Lambda,w\in\N$ be such that every integer $1\le t\le w$ divides $\Lambda$. Then $(j+k+1)!\mid\Lambda\,j!\,k!$ whenever $j+k+1\le w$.
Consequently, for every prime $\ell$ and every $w$, the integer $m=\prod_{s=1}^{w}s'$, where $s'$ is the prime-to-$\ell$ part of $s$, is prime
to $\ell$, and $\Lambda=m\,\ell^{\lfloor\log_\ell w\rfloor}$ satisfies $(j+k+1)!\mid\Lambda\,j!\,k!$ for $j+k+1\le w$.
\formallabel{GlobalTheory.T.beta_integral_denominator_bound}
\formallabel{GlobalTheory.LeviCut.betaDen}
\formallabel{GlobalTheory.BetaDenProof.beta_dvd}\formallabel{GlobalTheory.BetaDenProof.dvd_lam_of_le}
\end{lemma}

\begin{proof}
We prove by induction on $k$ that $(j+k+1)!\mid\Lambda\,j!\,k!$ for all $j$ with $j+k+1\le w$. For $k=0$ this is $(j+1)!\mid\Lambda\,j!$, which
holds since $j+1\mid\Lambda$. For the step, let $j+k+2\le w$. By induction (for $j$ and for $j+1$),
$(j+k+1)!\mid\Lambda\,j!\,k!$ and $(j+k+2)!\mid\Lambda\,(j+1)!\,k!$. The identity
\[
\Lambda\,j!\,(k+1)!+\Lambda\,(j+1)!\,k!=(j+k+2)\,\Lambda\,j!\,k!
\]
shows that $(j+k+2)!=(j+k+2)(j+k+1)!$ divides the left side, and since it divides the second summand, it divides
$\Lambda\,j!\,(k+1)!$. (This is the integrality of $\Lambda$ times the beta integral $B(j+1,k+1)=j!\,k!/(j+k+1)!$.)
For the second assertion: no factor $s'$ is divisible by $\ell$, so neither is $m$; and every $1\le t\le w$ is $t=\ell^{v}t'$ with $t'\mid m$ and
$\ell^v\le t\le w$, i.e.\ $v\le\lfloor\log_\ell w\rfloor$, so $t\mid\Lambda$.
\end{proof}

\begin{lemma}[The cut inclusion]\label{pf:sh-div:cutincl}
Let $\nu$ be antitone of length $p+q$, $c_1,c_2$ as in \ref{pf:sh-div:W:levi}, $K,w$ with $\nu_i\le\nu_j+K$ for all $i,j$ and $K+p+q\le w+1$,
and $\Lambda$ as in \ref{pf:sh-div:W:int}. Then $c_1c_2\Lambda^{pq}\nabla_\Z(\nu)\subseteq\Delta_\Z(\nu)$.
\formallabel{GlobalTheory.LeviCut.cut_incl}
\end{lemma}

\begin{proof}
Let $T$ be the endomorphism of \ref{pf:sh-div:W:levi} and $S_0=\mathrm{grid}_K(T)$. By \ref{pf:sh-div:W:gl}--\ref{pf:sh-div:W:adj} the operators
$D_{cd}=\mathrm{ad}(e_{cd})$ satisfy the hypotheses of \ref{pf:sh-div:W:inv}, and $D_{cd}T=e_{cd}T-Te_{cd}=0$ for $c,d$ in the same block;
hence $S_0$ commutes with all $e_{cd}$. The highest weight vector $v_\nu$ has weight $\nu$, is killed by the raising operators $e_{ij}$
($i<j$), and $\phi$ vanishes on the image of the lowering operators; so \ref{pf:sh-div:W:sc} gives $\phi(S_0v_\nu)z=\phi(Tv_\nu)=\phi(v_\nu)$ with
$z=\prod_{i<p\le j<p+q}(\nu_i-\nu_j+j-i)$, a non-zero integer. By \ref{pf:sh-div:W:schur} applied to $zS_0$, $zS_0=1$. Now let
$x\in\nabla_\Z(\nu)$. By \ref{pf:sh-div:W:int} (with $L=\nabla_\Z(\nu)$, $M=\Delta_\Z(\nu)$, $c=c_1c_2$), $c_1c_2\Lambda^{pq}S_0x\in\Delta_\Z(\nu)$, hence
$c_1c_2\Lambda^{pq}x=c_1c_2\Lambda^{pq}\,zS_0x=z\cdot(c_1c_2\Lambda^{pq}S_0x)\in\Delta_\Z(\nu)$.
\end{proof}

\begin{theorem}[The cut inequality]\label{pf:sh-div:excut}
$\sfname{ExCut}$ holds.
\formallabel{GlobalTheory.LeviCut.exCut}
\formallabel{GlobalTheory.LeviCut.ExCut₂}
\end{theorem}

\begin{proof}
Let $N,\lambda,\ell,\mathrm{lo},w,a$ be as in $\sfname{ExCut}$ and $b=N-a$, $\nu=\lambda|_N$. The blocks $\lambda|_a$ and $\lambda(a+\cdot)|_b$ are
antitone, and by the definition of $\shdivex_\ell$ there are $m_1,m_2$ prime to $\ell$ with $m_1\ell^{\shdivex_\ell(\lambda|_a)}\nabla_\Z(\lambda|_a)\subseteq
\Delta_\Z(\lambda|_a)$ and $m_2\ell^{\shdivex_\ell(\lambda(a+\cdot)|_b)}\nabla_\Z(\lambda(a+\cdot)|_b)\subseteq\Delta_\Z(\lambda(a+\cdot)|_b)$
(\lref{Formalism.CohomologicalDescent.LA.exact_z_specialization}). By \ref{pf:sh-div:W:beta} there is $m$ prime to $\ell$ such that
$\Lambda=m\ell^{\lfloor\log_\ell w\rfloor}$ satisfies the divisibility of \ref{pf:sh-div:W:int}. From $\mathrm{lo}\le\lambda(i)+N-1-i\le\mathrm{lo}+w$ we get
$\lambda(i)+j\le\lambda(j)+w+i$ for $i\le j<N$; with $i=0$, $j=N-1$ and antitonicity this gives $\nu_i\le\nu_j+K$ for all $i,j$, where
$K=w+1-N$, and $K+N\le w+1$. \Cref{pf:sh-div:cutincl} with $c_1=m_1\ell^{\shdivex_\ell(\lambda|_a)}$, $c_2=m_2\ell^{\shdivex_\ell(\lambda(a+\cdot)|_b)}$
and $\Lambda$ gives $m'\ell^{e}\nabla_\Z(\nu)\subseteq\Delta_\Z(\nu)$ with $m'=m_1m_2m^{ab}$ prime to $\ell$ and
$e=\shdivex_\ell(\lambda|_a)+\shdivex_\ell(\lambda(a+\cdot)|_b)+ab\lfloor\log_\ell w\rfloor$. By minimality, $\shdivex_\ell(\nu)\le e$.
\end{proof}

\begin{remark}\label{pf:sh-div:rem:excut}
The formal statement $\sfname{ExCut}$ is the special case $\GL_a\times\GL_{N-a}\subseteq\GL_N$ of a sub-additivity of the $\ell$-adic
exponents of the inclusion $\Delta(\nu)\subseteq\nabla(\nu)$ of integral Weyl modules \cite{Jantzen2003} along Levi subgroups; we do not know
a published reference for it in this form. Of the nine statements \ref{pf:sh-div:W:levi}--\ref{pf:sh-div:W:beta}, which are proved in the development, only
\ref{pf:sh-div:W:schur} and \ref{pf:sh-div:W:beta} are proved here (\cref{pf:sh-div:schur,pf:sh-div:betaden}).
\end{remark}
